\chapter{同伦论}
\label{cha:homotopy}

\index{acceptance|(}

在本章中，我们将在类型论内部发展一些同伦论。
我们采用 \cref{cha:basics} 中引入的同伦论的\emph{综合方法}：
空间、点、道路和同伦是基本概念，它们由类型及类型的元素来表示，特别是恒等类型。
道路与同伦的代数结构则由类型上自然的 $\infty$-群胚
\index{.infinity-groupoid@$\infty$-groupoid}%
结构所表示，这一结构由恒等类型的规则生成。
利用 \cref{cha:hits} 中引入的高阶归纳类型，我们可以通过万有性质直接描述空间。

\index{synthetic mathematics}%
这种综合方法有几个有趣的方面。
首先，它结合了具体模型（如拓扑空间\index{topological!space}或单纯集\index{simplicial!sets}）的优势与同伦论抽象范畴框架（如 Quillen 模型范畴\index{Quillen model category}）的优势。
一方面，我们的证明直观上很初等，具体地涉及类型中的点、道路和同伦。
另一方面，我们的方法又并不依赖于这些对象的任何具体呈现方式——例如，道路拼接的结合律是通过道路归纳证明的，而不是通过对映射 $[0,1] \to X$ 的重参数化或通过 horn 填充条件来证明。
类型论可以说是研究 $\infty$-群胚的抽象同伦论的一种非常便利的方式：
通过运用恒等类型的规则，我们可以避开许多 $\infty$-群胚定义中涉及的复杂组合学，并且在任何特定的证明中只需阐明所需的那部分结构。

类型论的抽象性意味着我们的证明能自动适用于多种不同的设定。
具体而言，如前所述，同伦类型论在 Kan\index{Kan complex} 单纯集\index{simplicial!sets}中有一种解释，而 Kan 单纯集是 $\infty$-群胚同伦论的一个模型。
因此，我们的证明适用于这个模型，并且通过沿着从单纯集到拓扑空间的几何实现\index{geometric realization}函子做转移，就能得到经典同伦论中相应定理的证明。
然而，尽管细节方面的工作仍在进行中，我们还可以在多种其他类似于 $\infty$-群胚范畴的范畴中解释类型论，例如 $(\infty,1)$-意象\index{.infinity1-topos@$(\infty,1)$-topos}。
因此，在类型论中证明一个结果，也就表明它在这些设定中同样成立。
这种额外的普适性是普通范畴逻辑的已知特性：泛等基础将它也扩展到了同伦论。

其次，我们的综合方法还启发了一些新的类型论方法与证明。
我们的一些证明是经典证明相当直接的转写。
另一些则更具类型论的特色，主要由 $\infty$-群胚运算的计算构成，其风格与计算机科学家用类型论推理计算机程序的方式非常相似。
使这些新证明成为可能的一个原因似乎是，类型论强调了同伦论中不同于其他方法的侧面：
尽管在 Kan\index{Kan complex} 单纯集这样的设定中，像道路归纳和高阶归纳类型的万有性质这样的工具也是可用的，但类型论提升了它们的重要性，因为它们是可用于处理这些类型的\emph{唯一}原始工具。
专注于这些工具，已经导出了对诸如圆周的万有覆叠和 Hopf 纤维化等熟悉构造的新描述，而这些描述仅仅使用了高阶归纳类型的递归原理。
这些描述非常直接，本章的许多证明都涉及对此类纤维化的计算。
\index{mathematics!constructive}%
我们证明的另一个新特点是它们是构造性的（假设泛等性与高阶归纳类型是构造性的）；
我们将在 \cref{sec:moreresults} 中描述它对于球面同伦群的一个应用。

\index{mathematics!formalized}%
\indexsee{formalization of mathematics}{mathematics, formalized}%
第三，我们的综合方法非常适合利用 \Coq 和 \Agda 等\index{proof!assistant}证明助理进行计算机验证。
本章所描述的证明，几乎全部都已通过计算机验证；其中许多证明最初是在证明助理中给出的，而后才为本书``去形式化''。
计算机验证证明的篇幅和工作量与这里展示的非形式证明不相上下，有时甚至更短、更容易完成。

\mentalpause

在着手展示我们的结果之前，我们为不熟悉同伦论的读者简要复习一些基本概念和定理。
同时，我们也概览一下本章将要证明的结果。

\index{classical!homotopy theory|(}%
同伦论是代数拓扑的一个分支，运用抽象代数的工具（如群论）来研究空间的性质。
同伦论研究者探讨的问题之一是如何判断两个空间是否``相同''，这里的``相同''指的是\emph{同伦等价}
\index{homotopy!equivalence!topological}%
（存在来回的连续映射能在同伦的意义下复合为恒等映射——这为``修正''那些不能严格复合为恒等的映射提供了可能）。
判断两个空间是否相同的常用方法，是计算与空间相关的\emph{代数不变量}，其中包括其\emph{同伦群}、\emph{同调群}与\emph{上同调群}。
\index{homology}%
\index{cohomology}%
等价的空间具有同构的同伦群和（上）同调群，因此，若两个空间的这些群不同，它们便不等价。
于是，这些代数不变量提供了关于空间的整体信息，可以用来区分空间，与连续性等概念所提供的局部信息形成互补。
例如，环面局部上看与 $2$-球面相似，但它有一个整体上的差异，因为它中间有一个洞，而这个差异在这两个空间的同伦群中是可见的。

同伦群最简单的例子是空间的\emph{基本群}
\index{fundamental!group}%
，记为 $\pi_1(X,x_0$)：给定空间 $X$ 及其中的一个点 $x_0$，可以构成一个群，其元素是 $x_0$ 处的环路（即从 $x_0$ 到 $x_0$ 的连续道路），并在同伦的意义下视作相同；群的运算由恒等道路（静止不动）、道路拼接和道路反转给出。
例如，$2$-球面的基本群是平凡的，而环面的基本群则不然，这表明球面和环面不是同伦等价的。
直观上可以这样理解：球面上的每一个环路都同伦于恒等道路，因为它内部可以被``填满''。
相比之下，环面上穿过甜甜圈那个孔洞的环路则无法同伦于恒等道路，因此其基本群中存在非平凡的元素。

\index{homotopy!group}
\emph{高阶同伦群}提供了关于空间的更多信息。
固定空间 $X$ 中的一点 $x_0$，考虑恒等道路 $\refl{x_0}$。
那么，$\refl{x_0}$ 与自身之间的同伦的同伦类构成一个群 $\pi_2(X,x_0)$，它揭示了空间某些二维结构的信息。
进而，$\pi_3(X,x_0$) 是同伦之间的同伦的同伦类构成的群，依此类推。
代数拓扑的基本问题之一就是\emph{计算空间 $X$ 的同伦群}，这意味着要给出 $\pi_k(X,x_0)$ 与某个更直接的群描述（例如通过乘法表或群表示给出的群）之间的群同构。
出人意料的是，这是个非常困难的问题，即使对球面这样简单的空间也是如此。
正如从 \cref{tab:homotopy-groups-of-spheres} 表中可以看到的那样，球面的高阶同伦群中会出现一些模式，但目前尚不存在通用公式，并且许多球面同伦群至今仍是未知的。

\bgroup
% Colors for the table of homotopy groups of spheres
\definecolor{xA}{\OPTcolormodel}{\OPTcolxA}
\definecolor{xB}{\OPTcolormodel}{\OPTcolxB}
\definecolor{xC}{\OPTcolormodel}{\OPTcolxC}
\definecolor{xD}{\OPTcolormodel}{\OPTcolxD}
\definecolor{xE}{\OPTcolormodel}{\OPTcolxE}
\definecolor{xF}{\OPTcolormodel}{\OPTcolxF}
\definecolor{xG}{\OPTcolormodel}{\OPTcolxG}
\definecolor{xH}{\OPTcolormodel}{\OPTcolxH}
\definecolor{xI}{\OPTcolormodel}{\OPTcolxI}
\definecolor{xJ}{\OPTcolormodel}{\OPTcolxJ}
\definecolor{xK}{\OPTcolormodel}{\OPTcolxK}
\definecolor{xL}{\OPTcolormodel}{\OPTcolxL}
\definecolor{xM}{\OPTcolormodel}{\OPTcolxM}

\newcommand{\cA}{\colorbox{xG}{\hbox to 20pt {\hfil$0$\hfil}}}
\newcommand{\cB}{\colorbox{xH}{\hbox to 20pt {\hfil$\Z$\hfil}}}
\newcommand{\cC}{\colorbox{xI}{\hbox to 20pt {\hfil$\Z_{2}$\hfil}}}
\newcommand{\cD}{\colorbox{xJ}{\hbox to 20pt {\hfil$\Z_{2}$\hfil}}}
\newcommand{\cE}{\colorbox{xK}{\hbox to 20pt {\hfil$\Z_{24}$\hfil}}}
\newcommand{\cF}{\colorbox{xL}{\hbox to 20pt {\hfil$0$\hfil}}}
\newcommand{\cG}{\colorbox{xM}{\hbox to 20pt {\hfil$0$\hfil}}}
\newcommand{\cH}{\colorbox{xF}{\hbox to 20pt {\hfil$0$\hfil}}}
\newcommand{\cI}{\colorbox{xE}{\hbox to 20pt {\hfil$0$\hfil}}}
\newcommand{\cJ}{\colorbox{xD}{\hbox to 20pt {\hfil$0$\hfil}}}
\newcommand{\cK}{\colorbox{xC}{\hbox to 20pt {\hfil$0$\hfil}}}
\newcommand{\cL}{\colorbox{xB}{\hbox to 20pt {\hfil$0$\hfil}}}
\newcommand{\cM}{\colorbox{xA}{\hbox to 20pt {\hfil$0$\hfil}}}

\begin{table}[htb]
\centering\small
\begin{tabular}{p{15pt}>{\centering\arraybackslash}p{\OPTspherescolwidth}>{\centering\arraybackslash}p{\OPTspherescolwidth}>{\centering\arraybackslash}p{\OPTspherescolwidth}>{\centering\arraybackslash}p{\OPTspherescolwidth}>{\centering\arraybackslash}p{\OPTspherescolwidth}>{\centering\arraybackslash}p{\OPTspherescolwidth}>{\centering\arraybackslash}p{\OPTspherescolwidth}>{\centering\arraybackslash}p{\OPTspherescolwidth}>{\centering\arraybackslash}p{\OPTspherescolwidth}}
\toprule
           & $\Sn^0$ & $\Sn^1$ & $\Sn^{2}$ & $\Sn^{3}$ & $\Sn^{4}$ & $\Sn^{5}$ & $\Sn^{6}$ & $\Sn^{7}$ & $\Sn^{8}$ \\ \addlinespace[3pt] \midrule
$\pi_{1}$  & $0$     & $\Z$    & \cA       & \cH       & \cI       & \cJ       & \cK       & \cL       & \cM       \\ \addlinespace[3pt]
$\pi_{2}$  & $0$     & $0$     & \cB       & \cA       & \cH       & \cI       & \cJ       & \cK       & \cL       \\ \addlinespace[3pt]
$\pi_{3}$  & $0$     & $0$     & $\Z$      & \cB       & \cA       & \cH       & \cI       & \cJ       & \cK       \\ \addlinespace[3pt]
$\pi_{4}$  & $0$     & $0$     & $\Z_{2}$  & \cC       & \cB       & \cA       & \cH       & \cI       & \cJ       \\ \addlinespace[3pt]
$\pi_{5}$  & $0$     & $0$     & $\Z_{2}$  & $\Z_{2}$  & \cC       & \cB       & \cA       & \cH       & \cI       \\ \addlinespace[3pt]
$\pi_{6}$  & $0$     & $0$     & $\Z_{12}$ & $\Z_{12}$ & \cD       & \cC       & \cB       & \cA       & \cH       \\ \addlinespace[3pt]
$\pi_{7}$  & $0$     & $0$     & $\Z_{2}$  & $\Z_{2}$  & {\footnotesize $\Z {\times} \Z_{12}$} & \cD & \cC & \cB     & \cA    \\ \addlinespace[3pt]
$\pi_{8}$  & $0$     & $0$     & $\Z_{2}$  & $\Z_{2}$  & $\Z_{2}^{2}$ & \cE & \cD & \cC & \cB \\ \addlinespace[3pt]
$\pi_{9}$  & $0$     & $0$     & $\Z_{3}$  & $\Z_{3}$  & $\Z_{2}^{2}$ & $\Z_{2}$ & \cE & \cD & \cC \\ \addlinespace[3pt]
$\pi_{10}$ & $0$     & $0$     & $\Z_{15}$ & $\Z_{15}$ & \footnotesize{$\Z_{24} {\times} \Z_{3}$} & $\Z_{2}$ & \cF & \cE & \cD \\ \addlinespace[3pt]
$\pi_{11}$ & $0$     & $0$     & $\Z_{2}$  & $\Z_{2}$  & $\Z_{15}$ & $\Z_{2}$ & $\Z$ & \cF & \cE \\ \addlinespace[3pt]
$\pi_{12}$ & $0$     & $0$     & $\Z_{2}^{2}$ & $\Z_{2}^{2}$ & $\Z_{2}$ & $\Z_{30}$ & $\Z_{2}$ & \cG & \cF \\ \addlinespace[3pt]
$\pi_{13}$ & $0$     & $0$     & {\footnotesize $\Z_{12} {\times} \Z_{2}$} & {\footnotesize $\Z_{12} {\times} \Z_{2}$} & $\Z_{2}^{3}$ & $\Z_{2}$ & $\Z_{60}$ & $\Z_{2}$ & \cG \\ \addlinespace[3pt]
%$\pi_{14}$ & $0$     & $0$     & $\Z_{84} {\times} \Z_{2}^{2}$ & $\Z_{84} {\times} \Z_{2}^{2}$ & {\footnotesize $\Z_{120} {\times} \Z_{12} {\times} \Z_{2}$} & $\Z_{2}^{3}$ & $\Z_{24} {\times} \Z_{2}$ & $\Z_{120}$ & $\Z_{2}$ \\ \addlinespace[3pt]
%$\pi_{15}$ & $0$     & $0$     & $\Z_{2}^{2}$ & $\Z_{2}^{2}$ & $\Z_{84} {\times} \Z_{2}^{5}$ & $\Z_{72} {\times} \Z_{2}$ & $\Z_{2}^{3}$ & $\Z_{2}^{3}$ & $\Z {\times} \Z_{120}$ \\ \addlinespace[3pt]
\bottomrule
\end{tabular}

\caption{球面的同伦群~\cite{wikipedia-groups}。\index{homotopy!group!of sphere}
表中每一格所列的群即为 $n$ 维球面 $\Sn^n$ 的第 $k$ 个同伦群 $\pi_k$ 所同构的群；其中 $\Z$ 是整数加群，
\index{integers}$\Z_{m}$ 是 $m$ 阶循环群。\index{group!cyclic}\index{cyclic group}
}
\label{tab:homotopy-groups-of-spheres}
\end{table}


\egroup

理解这种复杂性的一种途径，是借助 \cref{cha:basics} 中引入的空间与 $\infty$-群胚（$\infty$-groupoid）之间的对应关系。
\index{.infinity-groupoid@$\infty$-groupoid}%
如 \cref{sec:circle} 所述，2-球面由只包含一个点和一条 2 维环路的高阶归纳类型（higher inductive type）给出。
因此，我们可能会问：既然类型 $\Sn ^2$ 没有任何生成元能产生 3 维单元，为什么 $\pi_3(\Sn ^2)$ 会是 $\Z$ 呢？
原来，$\pi_3(\Sn ^2)$ 的生成元是利用 \cref{thm:EckmannHilton} 的证明中所描述的互换律构造出来的：
$\infty$-群胚的代数结构包含不同层级之间的非平凡交互，而这些交互便产生了更高维同伦群的元素。

\index{classical!homotopy theory|)}%

类型论为研究这一结构提供了自然的框架，因为我们可以方便地定义高阶同伦群。
回忆 \cref{def:loopspace}，对于 $n:\N$，带点类型 $(A,a)$ 的 $n$ 重迭代环路空间\index{loop space!iterated}递归地定义为：
\begin{align*}
  \Omega^0(A,a)&=(A,a)\\
  \Omega^{n+1}(A,a)&=\Omega^n(\Omega(A,a)).
\end{align*}
%
由此我们得到 $n$ 维环路的\emph{空间}（即一个类型），而这个空间本身又含有更高阶的同伦结构。\index{loop!n-@$n$-}
通过截断，我们可以得到 $n$ 维环路的集合（这一构造也在 \cref{sec:free-algebras} 中作为例子定义过）：

\begin{defn}[同伦群]\label{def-of-homotopy-groups}
  给定 $n\ge 1$ 和一个带点类型 $(A,a)$，我们定义 $A$ 在 $a$ 处的 \define{同伦群（homotopy groups）}
  \indexdef{homotopy!group}%
  为：
  \[\pi_n(A,a)\defeq \Trunc0{\Omega^n(A,a)}\]
\end{defn}

\noindent
由于 $n\ge 1$，$\Omega^n(A)$ 上的道路拼接与取逆运算诱导出 $\pi_n(A)$ 上的运算，从而自然地得到群结构。
若 $n\ge 2$，则群 $\pi_n(A)$ 是交换群\index{group!abelian}，这由 Eckmann--Hilton 论证\index{Eckmann--Hilton argument}（\cref{thm:EckmannHilton}）可知。
为了方便，我们也可以记 $\pi_0(A) \defeq \trunc0 A$，
但这个情形有所不同：它不仅不是群，而且在定义时也不依赖于 $A$ 中的任何基点。

\index{presentation!of an infinity-groupoid@of an $\infty$-groupoid}%
这个定义适合研究同伦群，因为类型 $X$ 的（高阶）归纳定义将 $X$ 呈现为一个自由类型，类似于一个自由 $\infty$-群胚，
\index{.infinity-groupoid@$\infty$-groupoid}%
而这种呈现\emph{决定了}但并不\emph{显式描述} $X$ 的高阶相等类型。
相等类型中既包含生成元（例如对于圆周来说是 $\lloop$），也包含对这些生成元施加所有群胚运算（恒等、复合、取逆、结合性、互换、\ldots）所得的结果。
因此，空间的高阶归纳呈现允许我们提出``$X$ 的相等类型实际上是什么？''这样的问题，尽管回答它可能需要相当多的数学论证。
这是类型论中一个熟知事实的高维推广：刻画 $X$ 的相等类型可能需要费一番功夫，即使 $X$ 只是一个普通的归纳类型，比如自然数或布尔值。
例如，$\bfalse$ 与 $\btrue$ 不相等这一定理并不能直接从定义推出；
参见 \cref{sec:compute-coprod}。

\index{univalence axiom}%
泛等公理在同伦群的计算中起着必不可少的作用（若无泛等公理，类型论可以有一种解释，使得所有路径——例如圆周上的环路——都是自反性路径）。这一点在下面计算圆周基本群时就会看到：从 $\Omega(\Sn^1)$ 到 $\Z$ 的映射是这样定义的：先将圆周上的一个环路映为集合 $\Z$ 的一个\index{automorphism!of Z, successor@of $\Z$, successor}自同构，使得例如 $\lloop \ct \opp \lloop$ 被映为 $\mathsf{successor} \ct \mathsf{predecessor}$（这里 $\mathsf{successor}$ 与 $\mathsf{predecessor}$ 是 $\Z$ 的自同构，借助泛等性可被视作宇宙中的路径），然后再将该自同构作用于 $0$。泛等性在宇宙中产生非平凡的道路，我们正是利用这一点从高阶归纳类型中的道路提取信息。

在本章中，我们首先计算球面的一些同伦群，包括 $\pi_k(\Sn ^1)$ (\cref{sec:pi1-s1-intro})、$k<n$ 时的 $\pi_{k}(\Sn ^n)$ (\cref{sec:conn-susp,sec:pik-le-n})、借助 Hopf 纤维化 (\cref{sec:hopf}) 与长正合列论证 (\cref{sec:long-exact-sequence-homotopy-groups}) 计算的 $\pi_2(\Sn ^2)$ 与 $\pi_3(\Sn ^2)$，以及借助 Freudenthal 纬悬定理 (\cref{sec:freudenthal}) 计算的 $\pi_n(\Sn ^n)$。接着，我们讨论 van Kampen 定理 (\cref{sec:van-kampen})——它刻画了推出的基本群——以及 Whitehead（怀特海德）原理的适用情况（即：一个在所有同伦群上诱导出等价的映射，何时本身也是等价）(\cref{sec:whitehead})。最后，我们简要概述一些未收入本书的进一步结果，例如 $n\ge 3$ 时的 $\pi_{n+1}(\Sn ^n)$、Blakers--Massey 定理，以及 Eilenberg（艾伦伯格）--Mac Lane（麦克兰恩）空间的一种构造 (\cref{sec:moreresults})。本章的预备知识包括 \cref{cha:typetheory,cha:basics,cha:hits,cha:hlevels} 以及 \cref{cha:logic} 的部分内容。

\index{fundamental!group!of circle|(}
\section{\texorpdfstring{$\pi_1(S^1)$}{π₁(S¹)}}
\label{sec:pi1-s1-intro}

本节的目标是证明 $\id {\pi_1(\Sn ^1)} {\Z}$。
事实上，我们将证明环路空间\index{loop space} ${\Omega(\Sn ^1)}$ 等价于 $\Z$。
这是一个更强的结论，因为根据定义 $\id{\pi_1(\Sn ^1)} {\trunc 0 {\Omega(\Sn ^1)}}$；因此，若 $\id {\Omega(\Sn ^1)} {\Z}$，则由合同性可得 $\id {\trunc
  0 {\Omega(\Sn ^1)}} {\trunc 0 {\Z}}$，而 $\Z$ 按定义是一个集合（因为它是一个集合商，参见 \cref{defn-Z,Z-quotient-by-canonical-representatives}），于是 $\id {\trunc 0 {\Z}} {\Z}$。
不仅如此，一旦知道 ${\Omega(\Sn ^1)}$ 是一个集合，便可推出当 $n>1$ 时 $\pi_n(\Sn^1)$ 是平凡的，这样一来，我们其实就计算出了 $\Sn^1$ 的\emph{全部}同伦群。

\subsection{开始探索}
\label{sec:pi1s1-initial-thoughts}

在 $\Omega(\Sn^1)$ 与 $\Z$ 之间定义两个方向的函数并不困难。
将 \cref{thm:looptothe} 特化到 $\lloop:\base=\base$ 的情形，便得到一个函数 $\lloop^{\blank} : \Z \rightarrow (\id{\base}{\base})$，其定义（粗略地说）为
\[
  \lloop^n =
  \begin{cases}
    \underbrace{\lloop \ct \lloop \ct \cdots \ct \lloop}_{n}  & \text{if $n > 0$,} \\
    \underbrace{\opp \lloop \ct \opp \lloop \ct \cdots \ct \opp \lloop}_{-n} & \text{if $n < 0$,} \\
    \refl{\base} & \text{if $n = 0$.}
\end{cases}
\]
%
而定义反方向的函数 $g:\Omega(\Sn^1)\to\Z$ 则要棘手一些。
注意到后继函数 $\Zsuc:\Z\to\Z$ 是一个等价，
\index{successor!isomorphism on Z@isomorphism on $\Z$}%
因而它在宇宙 $\type$ 中诱导出一条道路 $\ua(\Zsuc):\Z=\Z$。
于是，$\Sn^1$ 的递归原理诱导出一个映射 $c:\Sn^1\to\type$，由 $c(\base)\defeq \Z$ 与 $\apfunc c (\lloop) \defid \ua(\Zsuc)$ 给出。
这样我们就得到了 $\apfunc{c} : (\base=\base) \to (\Z=\Z)$，进而可以定义 $g(p)\defeq \transfib{X\mapsto X}{\apfunc{c}(p)}{0}$。

利用这些定义，借助 $n$ 的归纳原理 \cref{thm:sign-induction}，我们甚至能对任意 $n:\Z$ 证明 $g(\lloop^n)=n$。
（稍后我们将证明一个更一般的结论。）
然而，另一个等式 $\lloop^{g(p)}=p$ 的证明则要困难得多。
最直接的尝试是道路归纳，但道路归纳并不适用于像 $p:(\base=\base)$ 这样\emph{两个端点都固定}的环路！
这需要一个新的思路，它既可以从经典同伦论的角度，也可以从类型论的角度来阐释。
我们先从前者讲起。

\subsection{经典证明}
\label{sec:pi1s1-classical-proof}

\index{classical!homotopy theory|(}%
在经典同伦论中，证明 $\pi_1(\Sn^1)=\Z$ 有一个标准方法，要用到万有覆叠空间。
我们的证明可以看作此证明的类型论版本，其中覆叠空间体现为纤维是集合的纤维化。
\index{fibration}%
\index{total!space}%
回顾一下，同伦论中空间 $B$ 上的\emph{纤维化}对应于类型论中的类型族 $B\to\type$。
\index{path!fibration}%
\index{fibration!of paths}%
特别地，对于点 $x_0:B$，类型族 $(x\mapsto (x_0=x))$ 对应于\emph{道路纤维化} $P_{x_0} B \to B$：$P_{x_0} B$ 中的点是 $B$ 中始于 $x_0$ 的道路，而到 $B$ 的映射选取的是这条道路的另一个端点。
这个全空间 $P_{x_0} B$ 是可缩的，因为我们可以将任一道路``收缩''回其起始端点 $x_0$ —— 我们之前已经见过这一事实的类型论版本，即 \cref{thm:contr-paths}。
此外，$x_0$ 上的纤维就是环路空间\index{loop space} $\Omega(B,x_0)$ —— 在类型论中，由环路空间的定义即可看出。

\begin{figure}\centering
  \begin{tikzpicture}[xscale=1.4,yscale=.6]
    \node (R) at (2,1) {$\mathbb{R}$};
    \node (S1) at (2,-2) {$S^1$};
    \draw[->] (R) -- node[auto] {$w$} (S1);
    \draw (0,-2) ellipse (1 and .4);
    \draw[dotted] (1,0) arc (0:-30:1 and .8);
    \draw (1,0) arc (0:90:1 and .8) arc (90:270:1 and .3) coordinate (t1);
    \draw[white,line width=4pt] (t1) arc (-90:90:1 and .8);
    \draw (t1) arc (-90:90:1 and .8) arc (90:270:1 and .3) coordinate (t2);
    \draw[white,line width=4pt] (t2) arc (-90:90:1 and .8);
    \draw (t2) arc (-90:90:1 and .8) arc (90:270:1 and .3) coordinate (t3);
    \draw[white,line width=4pt] (t3) arc (-90:90:1 and .8);
    \draw (t3) arc (-90:-30:1 and .8) coordinate (t4);
    \draw[dotted] (t4) arc (-30:0:1 and .8);
    \node[fill,circle,inner sep=1pt,label={below:\scriptsize \base}] at (0,-2.4) {};
    \node[fill,circle,inner sep=1pt,label={above left:\scriptsize 0}] at (0,.2) {};
    \node[fill,circle,inner sep=1pt,label={above left:\scriptsize 1}] at (0,1.2) {};
    \node[fill,circle,inner sep=1pt,label={above left:\scriptsize 2}] at (0,2.2) {};
  \end{tikzpicture}
  \caption{经典拓扑中的环绕映射}\label{fig:winding}
\end{figure}

现在，在经典同伦论里——那里将 $\Sn^1$ 视为一个拓扑空间——我们可以这样进行。
\index{winding!map}%
考虑``环绕''映射 $w:\mathbb{R}\to \Sn ^1$，它看起来像一条投影到圆上的螺旋线（见 \cref{fig:winding}）。
映射 $w$ 将螺旋线\index{helix}上的每个点映为它正下方的圆上的点。
它是一个纤维化，并且每个点上的纤维同构于整数。\index{integers}
如果我们将底部沿逆时针方向绕行环路的那条道路提升，我们在螺旋线上就上升一层，使纤维中的整数加一。
类似地，沿底部环路顺时针绕行，对应于在螺旋线上下降一层，使该数减一。
这个纤维化称为圆的\emph{万有覆叠}。
\index{universal!cover}%
\index{cover!universal}%
\index{covering space!universal}%

经典同伦论中的一个基本事实是：基空间 $B$ 上的两个纤维化 $E_1$ 与 $E_2$ 之间的一个映射，若是 $E_1$ 与 $E_2$ 之间的同伦等价，则它在所有纤维上都诱导出同伦等价。
（我们此前也已经见过这一事实的类型论版本，见 \cref{thm:total-fiber-equiv}。）
由于 $\mathbb{R}$ 和 $P_{\base} S^1$ 都是可缩的拓扑空间，故它们同伦等价；因此，它们在基点处的纤维 $\Z$ 与 $\Omega(\Sn ^1)$ 也同伦等价。

\index{classical!homotopy theory|)}%

\subsection{类型论中的万有覆叠}
\label{sec:pi1s1-universal-cover}

\index{universal!cover|(}%
\index{cover!universal|(}%
\index{covering space!universal|(}%

我们来考虑如何在类型论中表达前述证明。
我们已经提到，$\Sn^1$ 的道路纤维化由类型族 $(x\mapsto (\base=x))$ 表示。
对于 $\Sn^1$ 的万有覆叠，我们也已经有了一个很好的候选：它恰恰就是我们在 \cref{sec:pi1s1-initial-thoughts} 中定义的类型族 $c:\Sn^1\to\type$！
根据定义，这个族在 $\base$ 处的纤维是 $\Z$，而沿着 $\lloop$ 传输的效果是加一 —— 因此它的行为正合我们对 \cref{fig:winding} 的期望。

不过，由于我们尚不知道这个族是否具备万有覆叠应有的性质（例如，其全空间是否单连通），我们暂且给它换一个名字。
为便于查阅，我们在此重复其定义。
\index{integers}

\begin{defn}[$\Sn^1$ 的万有覆叠] \label{S1-universal-cover}
  通过圆周递归定义 $\code : \Sn ^1 \to \type$，满足
  \begin{align*}
    \code(\base) &\defeq \Z \\
    \apfunc{\code}({\lloop}) &\defid \ua(\Zsuc).
  \end{align*}
\end{defn}

我们在此简略强调一下这个类型族的定义，因为它与经典同伦论中通常定义覆叠空间的方式大不相同。
要通过圆周递归定义一个函数，我们需要在值域中找出一个点和一个环路。
在当前情形下，值域是 $\type$，我们选择的点是 $\Z$，这对应着我们的期望：万有覆叠的纤维应该是整数。
我们选择的环路是 $\Z$ 上的后继/前驱
\index{successor!isomorphism on Z@isomorphism on $\Z$}%
\index{predecessor!isomorphism on Z@isomorphism on $\Z$}%
同构，这对应着一个几何事实：在基空间中绕环路一圈，在螺旋线上就上升一层。
这部分证明需要用到泛等性，因为我们需要将 $\Z$ 上的一个\emph{非平凡}等价转化为一个相等。

我们称此纤维化为``码''的纤维化，因为它的元素是组合数据，充当圆周上道路的码：整数 $n$ 编码了绕圆 $n$ 圈的道路。

从这个定义出发，容易计算出：用 $\code$ 传输会将 $\lloop$ 映为后继函数，而将 $\opp{\lloop}$ 映为前驱函数：


\begin{lem} \label{lem:transport-s1-code}
\id{\transfib \code \lloop x} {x + 1} 且
\id{\transfib \code {\opp \lloop} x} {x - 1}。
\end{lem}

\begin{proof}
对于第一个等式，我们计算如下：
\begin{align}
{\transfib \code \lloop x}
&= \transfib {A \mapsto A} {(\ap{\code}{\lloop})} x \tag{by \cref{thm:transport-compose}}\\
&= \transfib {A \mapsto A} {\ua (\Zsuc)} x \tag{by computation for $\rec{\Sn^1}$}\\
&= x + 1 \tag{by computation for \ua}.
\end{align}
第二个等式由第一个推出，因为 $\transfib{B}{p}{\blank}$ 与 $\transfib{B}{\opp p}{\blank}$ 总是互逆的，故 $\transfib\code {\opp \lloop}{\blank}$ 必是 $\Zsuc$ 的逆。
\end{proof}

此刻我们就能看出最初的方法问题何在：我们只定义了纤维 $\Omega(\Sn^1)$ 和 \Z 上的 $f$ 与 $g$，但实际上应该定义整个在 $\Sn^1$ 上的\emph{纤维化的}态射。
在类型论中，这意味着我们本该定义具有类型
\begin{align}
  \prd{x:\Sn^1} &((\base=x) \to \code(x)) \qquad\text{and/or} \label{eq:pi1s1-encode}\\
  \prd{x:\Sn^1} &(\code(x) \to (\base=x))\label{eq:pi1s1-decode}
\end{align}
的函数，而不是只定义这些函数在 $x$ 为 \base 时的特殊情况。
这也印证了类型论中一个常见的观察：当试图证明某个归纳类型的特定实例具有某种性质时，通常更容易先将命题推广到该类型的\emph{所有}实例，然后再用归纳法证明。
从这个角度看，$\Omega(\Sn^1)=\Z$ 的证明与 \cref{sec:compute-coprod,sec:compute-nat} 中刻画余积及自然数的相等类型遵循着相同的模式。

至此，完成证明有两种途径。
我们可以继续模仿经典论证，构造~\eqref{eq:pi1s1-encode} 或~\eqref{eq:pi1s1-decode}（二者选一即可），证明全空间之间的同伦等价会诱导纤维上的等价，然后证明万有覆叠的全空间是可缩的。
第一个关于 $\Omega(\Sn^1)=\Z$ 的类型论证明就遵循了这个模式；我们称之为\emph{同伦论证明}。

然而后来我们发现，还存在另一种证明，它更具有类型论的味道，并且更贴近 \cref{sec:compute-coprod,sec:compute-nat} 中的证明风格。
在这一证明中，我们直接构造出~\eqref{eq:pi1s1-encode} 和~\eqref{eq:pi1s1-decode}，并通过计算证明它们互逆。
我们称此为\emph{编码--解码证明}，因为我们将函数~\eqref{eq:pi1s1-encode} 和~\eqref{eq:pi1s1-decode} 分别称为\emph{编码}与\emph{解码}。
两种证明都使用了前面给出的同一个覆叠构造。
前一种证明从全空间的等价诱导出纤维上的等价，而编码--解码证明则显式地构造出纤维之间的逆映射（即\emph{解码}）。
前一种证明使用了可缩性，而编码--解码证明使用了道路归纳、圆周归纳和整数归纳。
这些工具与证明可缩性所用的工具是一样的——事实上，道路归纳\emph{本质上就是}道路纤维化的可缩性再配上 $\mathsf{transport}$——只是应用方式不同。

既然这是一本关于同伦类型论的书，我们将首先介绍编码--解码证明。
若有同伦论背景的读者感到迷失，我们建议直接跳到同伦论证明（\cref{subsec:pi1s1-homotopy-theory}）。

\index{universal!cover|)}%
\index{cover!universal|)}%
\index{covering space!universal|)}%

\subsection{编码--解码证明}
\label{subsec:pi1s1-encode-decode}

\index{encode-decode method|(}%
\indexsee{encode}{encode-decode method}%
\indexsee{decode}{encode-decode method}%
我们从将道路映射到代码的函数~\eqref{eq:pi1s1-encode} 开始：


\begin{defn}
定义 $\encode : \prd{x : \Sn ^1} (\base=x) \rightarrow  \code(x)$ 为
\[
\encode \: p \defeq \transfib{\code} p 0
\]
（其中参数 $x$ 是隐式的）。
\end{defn}


编码的定义方式是将一条道路提升到万有覆叠，这决定了一个等价，然后将得到的等价应用到 $0$。
这个函数的有趣之处在于，当圆上的环路是用同伦类型论中抽象的群胚框架来表示时，它能从这样一个环路计算出一个具体的数。
为了直观理解它是如何做到的，我们可以根据前面的引理观察到：$\transfib \code \lloop x$ 是后继映射，而 $\transfib \code {\opp
  \lloop} x$ 是前驱映射。
此外，$\mathsf{transport}$ 具有函子性（\cref{cha:basics}），所以 $\transfib{\code} {\lloop \ct \lloop}{\blank}$ 就是
\[(\transfib \code \lloop-) \circ (\transfib \code \lloop-)\]
依此类推。
因此，当 $p$ 是一个像
\[
\lloop \ct \opp \lloop \ct \lloop \ct \cdots
\]
这样的复合时，$\transfib{\code}{p}{\blank}$ 将计算出函数的复合，例如
\[
\Zsuc \circ \Zpred \circ \Zsuc \circ \cdots
\]
将这一组函数的复合应用到 $0$，就会计算出该道路的\index{winding!number}%
\emph{环绕数}——即在消去逆元之后，它绕圆了多少圈，其方向由正负号标记。
由此可见，$\encode$ 的计算行为源于高阶归纳类型与泛等性的归约规则，以及 $\mathsf{transport}$ 对复合与逆的作用。

注意，特例 $\encode' \defeq \encode_{\base}$ 的类型是 $(\id \base \base) \rightarrow \Z$。
这将是我们所需等价的一半；实际上，它正是 \cref{sec:pi1s1-initial-thoughts} 中定义的函数 $g$。

类似地，函数~\eqref{eq:pi1s1-decode} 是 \cref{sec:pi1s1-initial-thoughts} 中函数 $\lloop^{\blank}$ 的推广。

\begin{defn}\label{thm:pi1s1-decode}
通过对 $x$ 作圆周归纳来定义 $\decode : \prd{x : \Sn ^1}\code(x) \rightarrow (\base=x)$。
我们只需给出一个函数
${\code(\base) \rightarrow (\base=\base)}$，为此我们使用 $\lloop^{\blank}$，并证明 $\lloop^{\blank}$ 尊重环路即可。
\end{defn}

\begin{proof}
为了证明 $\lloop^{\blank}$ 尊重环路，只需给出一条从 $\lloop^{\blank}$ 到其自身的道路，且这条道路位于 $\lloop$ 之上。
根据依值道路的定义，这意味着需要一条从
\[\transfib {(x' \mapsto \code(x') \rightarrow (\base=x'))} {\lloop} {\lloop^{\blank}}\]
到 $\lloop^{\blank}$ 的道路。
我们如下构造这样一条道路：
\begin{align*}
  \MoveEqLeft \transfib {(x' \mapsto \code(x') \rightarrow (\base=x'))} \lloop {\lloop^{\blank}}\\
&= \transfibf {x'\mapsto (\base=x')}(\lloop) \circ {\lloop^{\blank}} \circ \transfibf \code ({\opp \lloop}) \\
&= (- \ct \lloop) \circ (\lloop^{\blank}) \circ \transfibf \code ({\opp \lloop}) \\
&= (- \ct \lloop) \circ (\lloop^{\blank}) \circ \Zpred \\
&= (n \mapsto \lloop^{n - 1} \ct \lloop).
\end{align*}
在第一行，我们应用了当纤维化的外层连接词是 $\rightarrow$ 时 $\mathsf{transport}$ 的刻画，它将这个 $\mathsf{transport}$ 化归为在定义域和值域类型上与 $\mathsf{transport}$ 的前复合与后复合。
在第二行，我们应用了当类型族是 $x\mapsto \id{\base}{x}$ 时 $\mathsf{transport}$ 的刻画，即道路的后复合。
在第三行，我们使用了从 \cref{lem:transport-s1-code} 得到的 $\code$ 在 $\opp \lloop$ 上的作用。
在第四行，我们仅仅将函数复合化归了。
因此，只需证明对于所有 $n$，有 $\id{\lloop^{n - 1} \ct \lloop}{\lloop^{n}}$。
这是 \cref{thm:sign-induction} 的一个简单应用，利用了群胚律。
\end{proof}

现在我们可以证明 $\encode$ 和 $\decode$ 互为拟逆。
曾经困难的方向现在变得简单了！

\begin{lem} \label{lem:s1-decode-encode}
对于所有 $x: \Sn ^1$ 和 $p : \id \base x$，有 $\id
{\decode_x({{\encode_x(p)}})} p$。
\end{lem}

\begin{proof}
通过道路归纳，只需证明
\narrowequation{\id {\decode_{\base}({{\encode_{\base}(\refl{\base})}})} {\refl{\base}}.}
但是
\narrowequation{\encode_{\base}(\refl{\base}) \jdeq \transfib{\code}{\refl{\base}} 0 \jdeq 0,}
并且 $\decode_{\base}(0) \jdeq \lloop^ 0 \jdeq \refl{\base}$。
\end{proof}

另一个方向也不难多少。

\begin{lem} \label{lem:s1-encode-decode} 对于所有
$x: \Sn ^1$ 和 $c : \code(x)$，我们有 $\id
{\encode_x({{\decode_x(c)}})} c$。
\end{lem}

\begin{proof}
证明通过圆周归纳进行。只需证明 $\base$ 的情形，因为 $\lloop$ 的情形是 $\Z$ 中两条道路之间的道路，而 $\Z$ 是集合，所以这是直接可得的。

因此，我们需要证明：对于所有 $n : \Z$，有
\[
\id {\encode'(\lloop^n)} {n}.
\]
证明通过对 $n$ 的归纳进行，使用 \cref{thm:sign-induction}。
%
\begin{itemize}

\item 对于 $0$ 的情形，结果按定义成立。

\item 对于 $n+1$ 的情形，
\begin{align}
 {\encode'(\lloop^{n+1})}
&= {\encode'(\lloop^{n} \ct \lloop)} \tag{by definition of $\lloop^{\blank}$} \\
&= \transfib{\code}{(\lloop^{n} \ct \lloop)}{0} \tag{by definition of $\encode$}\\
&= \transfib{\code}{\lloop}{(\transfib{\code}{\lloop^n}{0})} \tag{by functoriality}\\
&= {(\transfib{\code}{\lloop^n}{0})} + 1 \tag{by \cref{lem:transport-s1-code}}\\
&= n + 1. \tag{by the inductive hypothesis}
\end{align}

\item 负数的情形是类似的。\qedhere
\end{itemize}
\end{proof}

最后，我们给出以下定理。

\begin{thm}
存在一族等价 $\prd{x : \Sn ^1} (\eqv {(\base=x)} {\code(x)})$。
\end{thm}

\begin{proof}
映射 $\encode$ 和 $\decode$ 互为拟逆，这由
\cref{lem:s1-decode-encode,lem:s1-encode-decode}
给出。
\end{proof}

在 \base 处实例化，得到


\begin{cor}\label{cor:omega-s1}
$\eqv {\Omega(\Sn^1,\base)} {\Z}$。
\end{cor}

一个简单的归纳表明该等价把加法对应到复合，因此有 $\Omega(\Sn ^1) = \Z$（作为群）。

\begin{cor} \label{cor:pi1s1}
$\id{\pi_1(\Sn ^1)} {\Z}$，而当 $n>1$ 时，$\id{\pi_n(\Sn ^1)}0$。
\end{cor}

\begin{proof}
对于 $n=1$，我们在上面从 \cref{cor:omega-s1} 概述了证明。
对于 $n > 1$，有 $\trunc 0 {\Omega^{n}(\Sn ^1)} = \trunc 0 {\Omega^{n-1}(\Omega{\Sn ^1})} = \trunc 0 {\Omega^{n-1}(\Z)}$。
由于 $\Z$ 是集合，$\Omega^{n-1}(\Z)$ 是可缩的，因此这是平凡的。
\end{proof}

\index{encode-decode method|)}%

\subsection{同伦论的证明}
\label{subsec:pi1s1-homotopy-theory}

在 \cref{sec:pi1s1-universal-cover} 中，我们在类型论里定义了假定的万有覆叠 $\code:\Sn^1\to\type$；在 \cref{subsec:pi1s1-encode-decode} 中，我们定义了一个从道路纤维化到万有覆叠的映射 $\encode : \prd{x:\Sn^1} (\base=x) \to \code(x)$。
对于经典证明，剩下的部分是证明：因为两个全空间都是可缩的，所以这个映射诱导了全空间上的等价，并由此推断它在每个纤维上也必然是等价。

\index{total!space}%
在 \cref{thm:contr-paths} 中我们看到全空间 $\sm{x:\Sn^1} (\base=x)$ 是可缩的。
对于另一个全空间，我们有：

\begin{lem}\label{thm:iscontr-s1cover}
  类型 $\sm{x:\Sn^1}\code(x)$ 是可缩的。
\end{lem}

\begin{proof}
  我们以如下取值应用展平引理（\cref{thm:flattening}）：
  \begin{itemize}
  \item $A\defeq\unit$ 且 $B\defeq\unit$，其中 $f$ 和 $g$ 取为显然的函数。
    因此，展平引理中的基础高阶归纳类型 $W$ 等价于 $\Sn^1$。
  \item $C:A\to\type$ 取常值 \Z。
  \item $D:\prd{b:B} (\eqv{\Z}{\Z})$ 取常值 $\Zsuc$。
  \end{itemize}
  于是，展平引理中定义的类型族 $P:\Sn^1\to\type$ 等价于 $\code:\Sn^1\to\type$。
  因此，展平引理告诉我们，$\sm{x:\Sn^1}\code(x)$ 等价于一个具有以下生成元的高阶归纳类型，我们将其记作 $R$：
  \begin{itemize}
  \item 一个函数 $\mathsf{c}: \Z \to R$。
  \item 对每个 $z:\Z$，一条道路 $\mathsf{p}_z:\mathsf{c}(z) = \mathsf{c}(\Zsuc(z))$。
  \end{itemize}
  我们可以将这个类型称为\define{同伦实数（homotopical reals）}；
  \indexdef{real numbers!homotopical}%
  它在经典证明中扮演的角色与拓扑空间 $\mathbb{R}$ 相同。

  因此，只需再证明 $R$ 是可缩的。
  我们选取 $\mathsf{c}(0)$ 作为收缩中心；现在需要证明对所有 $x:R$ 有 $x=\mathsf{c}(0)$。
  我们对 $R$ 进行归纳来证明。
  首先，当 $x$ 为 $\mathsf{c}(z)$ 时，我们需要给出一条道路 $q_z:\mathsf{c}(0) = \mathsf{c}(z)$，这可以通过对 $z:\Z$ 归纳来完成，使用 \cref{thm:sign-induction}：
  \begin{align*}
    q_0 &\defid \refl{\mathsf{c}(0)}\\
    q_{n+1} &\defid q_n \ct \mathsf{p}_n & &\text{for $n\ge 0$}\\
    q_{n-1} &\defid q_n \ct \opp{\mathsf{p}_{n-1}} & &\text{for $n\le 0$.}
  \end{align*}
  其次，我们需要证明，对任意 $z:\Z$，道路 $q_z$ 沿 $\mathsf{p}_z$ 传输后得到 $q_{z+1}$。
  根据道路传输的性质，这意味着我们需要 $q_z \ct \mathsf{p}_z = q_{z+1}$。
  利用 $q_z$ 的定义，通过对 $z$ 归纳容易完成。
  这就完成了 $R$ 可缩的证明，因而 $\sm{x:\Sn^1}\code(x)$ 也是可缩的。
\end{proof}

\begin{cor}\label{thm:encode-total-equiv}
  由 \encode 诱导的映射：
  \[ \tsm{x:\Sn^1} (\base=x) \to \tsm{x:\Sn^1}\code(x) \]
  是一个等价。
\end{cor}

\begin{proof}
  这两个类型都是可缩的。
\end{proof}

\begin{thm}
  $\eqv {\Omega(\Sn^1,\base)} {\Z}$。
\end{thm}

\begin{proof}
  对 $\encode$ 应用 \cref{thm:total-fiber-equiv}，并使用 \cref{thm:encode-total-equiv}。
\end{proof}

本质上，这两种证明并无太大不同：encode-decode 证明可以看作是同伦论证明的一种``约化''或``展开''。
二者各有优势；这两种观点之间的相互交织，也是本课题趣味的一部分。
\index{fundamental!group!of circle|)}

\subsection{将万有覆叠视作恒等系统}
\label{sec:pi1s1-idsys}

注意，纤维化 $\code:\Sn^1\to\type$ 连同 $0:\code(\base)$ 构成 \cref{defn:identity-systems} 意义下的一个\emph{带点谓词（pointed predicate）}。
从这个角度看，\cref{subsec:pi1s1-encode-decode} 中的编码-解码证明本质上是证明 \code 满足 \cref{thm:identity-systems}\ref{item:identity-systems3}，而 \cref{subsec:pi1s1-homotopy-theory} 中的同伦论证明则是证明它满足 \cref{thm:identity-systems}\ref{item:identity-systems4}。
这提示了第三种途径。

\begin{thm}
  在 \cref{defn:identity-systems} 的意义下，偶对 $(\code,0)$ 是 $\base:\Sn^1$ 处的一个恒等系统。
\end{thm}

\begin{proof}
  给定 $D:\prd{x:\Sn^1} \code(x) \to \type$ 及 $d:D(\base,0)$，我们希望定义函数 $f:\prd{x:\Sn^1}{c:\code(x)} D(x,c)$。
  由圆周归纳，只需指定 $f(\base):\prd{c:\code(\base)} D(\base,c)$ 并验证 $\trans{\lloop}{f(\base)} = f(\base)$。

  显然，$\code(\base)\jdeq \Z$。
  由 \cref{lem:transport-s1-code} 并对 $n$ 归纳，可以对任意整数 $n$ 得到一条道路 $p_n : \transfib{\code}{\lloop^n}{0} = n$。
  因此，由 $\Sigma$-类型中的道路，我们在 $\sm{x:\Sn^1} \code(x)$ 中得到一条道路 $\pairpath(\lloop^n,p_n) : (\base,0) = (\base,n)$。
  将 $d$ 沿此道路在 $D$ 相应的纤维化 $\widehat{D}:(\sm{x:\Sn^1} \code(x)) \to\type$ 中传输，便对任意 $n:\Z$ 得到 $D(\base,n)$ 中的一个元素。
  我们定义此元素为 $f(\base)(n)$：
  \[ f(\base)(n) \defeq \transfib{\widehat{D}}{\pairpath(\lloop^n,p_n)}{d}. \]
  %
  现在我们需要 $\transfib{\lam{x} \prd{c:\code(x)} D(x,c)}{\lloop}{f(\base)} = f(\base)$。
  由 \cref{thm:dpath-forall}，这意味着我们需要证明对任意 $n:\Z$，
  %
  \begin{narrowmultline*}
    \transfib{\widehat D}{\pairpath(\lloop,\refl{\trans\lloop n})}{f(\base)(n)}
    =_{D(\base,\trans\lloop n)} \narrowbreak
    f(\base)(\trans\lloop n).
  \end{narrowmultline*}
  %
  现在我们有一条道路 $q:\trans\lloop n = n+1$，故沿此道路传输后，只需证明
  \begin{multline*}
    \transfib{D(\base)}{q}{\transfib{\widehat D}{\pairpath(\lloop,\refl{\trans\lloop n})}{f(\base)(n)}}\\
    =_{D(\base,n+1)} \transfib{D(\base)}{q}{f(\base)(\trans\lloop n)}.
  \end{multline*}
  由关于 transport 与依值应用的几个引理，上式等价于
  \[ \transfib{\widehat D}{\pairpath(\lloop,q)}{f(\base)(n)} =_{D(\base,n+1)} f(\base)(n+1). \]
  然而，展开 $f(\base)$ 的定义，我们有
  \begin{narrowmultline*}
    \transfib{\widehat D}{\pairpath(\lloop,q)}{f(\base)(n)}
    \narrowbreak
    \begin{aligned}[t]
      &= \transfib{\widehat D}{\pairpath(\lloop,q)}{\transfib{\widehat{D}}{\pairpath(\lloop^n,p_n)}{d}}\\
      &= \transfib{\widehat D}{\pairpath(\lloop^n,p_n) \ct \pairpath(\lloop,q)}{d}\\
      &= \transfib{\widehat D}{\pairpath(\lloop^{n+1},p_{n+1})}{d}\\
      &= f(\base)(n+1).
    \end{aligned}
  \end{narrowmultline*}
  这里我们用到了 transport 的函子性、$\Sigma$-类型中复合的刻画（这是留给读者的练习）、以及一个将 $p_n$、$q$ 与 $p_{n+1}$ 联系起来的引理，该引理的陈述与证明也留给读者。

  这就完成了 $f:\prd{x:\Sn^1}{c:\code(x)} D(x,c)$ 的构造。
  由于
  \[f(\base,0) \jdeq \trans{\pairpath(\lloop^0,p_0)}{d} = \trans{\refl{\base}}{d} = d,\]
  我们已证明 $(\code,0)$ 是一个恒等系统。
\end{proof}

\begin{cor}
  对于任意 $x:\Sn^1$，有 $\eqv{(\base=x)}{\code(x)}$。
\end{cor}

\begin{proof}
  由 \cref{thm:identity-systems}。
\end{proof}

当然，这个证明本质上包含了与前两种证明相同的要素。
粗略地说，它统一了 \cref{thm:pi1s1-decode,lem:s1-encode-decode} 的证明，在一般情形下只需进行一次必要的归纳论证。

\begin{rmk}
  注意，以上所有关于 $\eqv{\pi_1(\Sn^1)}{\Z}$ 的证明都本质地使用了泛等公理。
  这是不可避免的：为了证明 $\eqv{\pi_1(\Sn^1)}{\Z}$，泛等性或其类似物是\emph{必需}的。
  在缺乏泛等性的情况下，可以一致地假设``所有类型都是集合''这一陈述（又称``恒等证明的唯一性''或``K 公理''，如 \cref{sec:hedberg} 中所讨论的），而这一陈述反而蕴含 $\eqv{\pi_1(\Sn^1)}{\unit}$。
  事实上，$\pi_1(\Sn^1)$ 的（非）平凡性恰好可以检测是否所有类型都是集合：\cref{thm:loop-nontrivial} 的证明反过来表明，如果 $\lloop=\refl{\base}$，那么所有类型都是集合。
\end{rmk}

\section{纬悬的连通性}
\label{sec:conn-susp}

回顾 \cref{sec:connectivity}，一个类型 $A$ 称为 \define{$n$-连通的（$n$-connected）}，如果 $\trunc nA$ 是可缩的。
本节的目标是证明 \cref{sec:suspension} 中定义的纬悬运算会提升连通度。

\begin{thm} \label{thm:suspension-increases-connectedness}
  如果 $A$ 是 $n$-连通的，那么 $A$ 的纬悬是 $(n+1)$-连通的。
\end{thm}

\begin{proof}
  我们在 \cref{sec:colimits} 中曾指出，$A$ 的纬悬是推出 $\unit\sqcup^A\unit$，因此需要证明下列类型是可缩的：
  %
  \[\trunc{n+1}{\unit\sqcup^A\unit}.\]
  %
  由 \cref{reflectcommutespushout} 可知
  $\trunc{n+1}{\unit\sqcup^A\unit}$ 是图表
  \[\xymatrix{\trunc{n+1}A \ar[d] \ar[r] & \trunc{n+1}{\unit} \\
    \trunc{n+1}{\unit} & }.\]
  %
  在 $\typelep{n+1}$ 中的推出。
  已知 $\trunc{n+1}{\unit}=\unit$，类型
  $\trunc{n+1}{\unit\sqcup^A\unit}$ 也是下列图表在 $\typelep{n+1}$ 中的推出（因为两个图表相等）：
  %
  \[\Ddiag=\vcenter{\xymatrix{\trunc{n+1}A \ar[d] \ar[r] & \unit \\
    \unit & }}.\]
  %
  我们现在证明 $\unit$ 也是 $\Ddiag$ 在 $\typelep{n+1}$ 中的推出。
  %
  令 $E$ 为一个 $(n+1)$-截断的类型；我们需要证明以下映射是一个等价：
  %
  \[\function{(\unit \to E)}{\cocone{\Ddiag}{E}}{y}
  {(y,y,\lamu{u:{\trunc{n+1}A}} \refl{y(\ttt)})}.\]
  %
  其中回顾 $\cocone{\Ddiag}{E}$ 是类型
  \[\sm{f:\unit \to E}{g:\unit \to E}(\trunc{n+1}A\to
  (f(\ttt)=_E{}g(\ttt))).\]
  %
  映射 $\function{(\unit\to E)}{E}{f}{f(\ttt)}$ 是一个等价，因此我们也有
  \[\cocone{\Ddiag}{E}=\sm{x:E}{y:E}(\trunc{n+1}A\to(x=_Ey)).\]
  %
  现在，$A$ 是 $n$-连通的，故 $\trunc{n+1}A$ 也是 $n$-连通的（因为
  $\trunc n{\trunc{n+1}A}=\trunc nA=\unit$），且 $(x=_Ey)$ 是 $n$-截断的（因为
  $E$ 是 $(n+1)$-截断的）。因此，由 \cref{connectedtotruncated}，以下映射是一个等价：
  %
  \[\function{(x=_Ey)}{(\trunc{n+1}A\to(x=_Ey))}{p}{\lam{z} p}\]
  %
  因此我们有
  %
  \[\cocone{\Ddiag}{E}=\sm{x:E}{y:E}(x=_Ey).\]
  %
  但下列映射是一个等价：
  %
  \[\function{E}{\sm{x:E}{y:E}(x=_Ey)}{x}{(x,x,\refl{x})}.\]
  %
  因此
  %
  \[\cocone{\Ddiag}{E}=E.\]
  %
  最终我们得到一个等价：
  %
  \[(\unit \to E)\eqvsym\cocone{\Ddiag}{E}\]
  %
  现在我们可以展开定义以得到该映射的显式表达式，并且容易看出这正是我们一开始的那个映射。

  因此我们证明了 $\unit$ 是 $\Ddiag$ 在 $\typelep{n+1}$ 中的推出。利用推出的唯一性，得到 $\trunc{n+1}{\unit\sqcup^A\unit}=\unit$，这证明了 $A$ 的纬悬是 $(n+1)$-连通的。
\end{proof}

\begin{cor} \label{cor:sn-connected}
  对所有 $n:\N$，球面 $\Sn^n$ 是 $(n-1)$-连通的。
\end{cor}

\begin{proof}
  我们对 $n$ 归纳证明。
  %
  对于 $n=0$，我们需要证明 $\Sn^0$ 仅知有实例，这是显然的。
  %
  设 $n:\N$ 且 $\Sn^n$ 是 $(n-1)$-连通的。根据定义，$\Sn^{n+1}$ 是 $\Sn^n$ 的纬悬，因此由前述引理可知 $\Sn^{n+1}$ 是 $n$-连通的。
\end{proof}

\section{$n$-连通空间的 \texorpdfstring{$\pi_{k \le n}$}{π\_(k≤n)} 与 \texorpdfstring{$\pi_{k < n}(\Sn ^n)$}{π\_(k<n)(Sⁿ)}}
\label{sec:pik-le-n}

令 $(A,a)$ 为一个带点类型且 $n:\N$。回顾
\cref{thm:homotopy-groups}，若 $n>0$，集合 $\pi_n(A,a)$ 具有群结构；若 $n>1$，该群是交换群\index{group!abelian}。

关于 $n$-截断类型与 $n$-连通类型的同伦群，我们现在可以有所讨论。

\begin{lem}
  如果 $A$ 是 $n$-截断的且 $a:A$，那么对所有 $k>n$ 有 $\pi_k(A,a)=\unit$。
\end{lem}

\begin{proof}
  $n$-类型的环路空间是 $(n-1)$-类型，因此 $\Omega^k(A,a)$ 是 $(n-k)$-类型，且 $(n-k)\le-1$，所以 $\Omega^k(A,a)$ 是一个命题。但 $\Omega^k(A,a)$ 是有实例的，故它实际上是可缩的，并且
  $\pi_k(A,a)=\trunc0{\Omega^k(A,a)}=\trunc0{\unit}=\unit$。
\end{proof}

\begin{lem} \label{lem:pik-nconnected}
  如果 $A$ 是 $n$-连通的且 $a:A$，那么对所有 $k\le{}n$ 有 $\pi_k(A,a)=\unit$。
\end{lem}

\begin{proof}
  我们有下列等式序列：
  %
  \begin{narrowmultline*}
    \pi_k(A,a) = \trunc0{\Omega^k(A,a)}
    = \Omega^k(\trunc k{(A,a)})
    = \Omega^k(\trunc k{\trunc n{(A,a)}})
    = \narrowbreak
      \Omega^k(\trunc k{\unit})
    = \Omega^k(\unit)
    = \unit.
  \end{narrowmultline*}
  %
  第三个等式利用了 $k\le{}n$ 这一事实，从而可以使用 $\truncf k\circ\truncf n=\truncf k$，而第四个等式利用了 $A$ 是 $n$-连通的这一事实。
\end{proof}

\begin{cor}
  当 $k < n$ 时，$\pi_k(\Sn ^n) = \unit$。
\end{cor}

\begin{proof}
  由 \cref{cor:sn-connected}，球面 $\Sn^n$ 是 $(n-1)$-连通的，因此我们可以应用 \cref{lem:pik-nconnected}。
\end{proof}

\section{纤维列与长正合列}
\label{sec:long-exact-sequence-homotopy-groups}

\index{fiber sequence|(}%
\index{sequence!fiber|(}%

如果一个函数 $f:X\to Y$ 的陪域配有一个基点 $y_0:Y$，那么 $f$ 在 $y_0$ 上的纤维 $F\defeq \hfib f {y_0}$ 就称为 \define{$f$ 的纤维 (the fiber of $f$)}\index{fiber}。
（若 $Y$ 是连通的，则 $F$ 仅在相差一个等价的意义下被确定；参见 \cref{ex:unique-fiber}。）
我们现在证明，如果 $X$ 也是带点类型且 $f$ 保持基点，那么 $F$、$X$ 和 $Y$ 的同伦群之间存在一个关系，具体表现为一个\emph{长正合列}。我们通过与这样一个 $f$ 相关联的\emph{纤维列}来推导这一结论。

\begin{defn}\label{def:pointedmap}
  带点类型 $(X,x_0)$ 与 $(Y,y_0)$ 之间的 \define{带点映射 (pointed map)}
  \indexdef{pointed!map}%
  \indexsee{function!pointed}{pointed map}%
  是一个映射 $f:X\to Y$，并附带一条道路 $f_0:f(x_0)=y_0$。
\end{defn}

对于任意带点类型 $(X,x_0)$ 和 $(Y,y_0)$，都存在一个带点映射 $(\lam{x} y_0) : X\to Y$，它是取值为基点的常值映射。
我们称此映射为 \define{零映射 (zero map)}\indexdef{zero!map}\indexdef{function!zero}，有时也将其记作 $0:X\to Y$。

回顾，每个带点类型 $(X,x_0)$ 都有一个环路空间 \index{loop space} $\Omega (X,x_0)$。
现在我们指出，这一运算在带点映射上具有函子性。\index{loop space!functoriality of}\index{functor!loop space}

\begin{defn}\label{def:loopfunctor}
  给定带点类型之间的带点映射 $f:X \to Y$，我们如下定义一个带点映射 $\Omega f:\Omega X \to \Omega Y$：
  \[(\Omega f)(p) \defeq \rev{f_0}\ct\ap{f}{p}\ct f_0.\]
  道路 $(\Omega f)_0 : (\Omega f) (\refl{x_0}) = \refl{y_0}$ 将 $\Omega f$ 展示为带点映射，它是类型如下的显然道路：
  \[\rev{f_0}\ct\ap{f}{\refl{x_0}}\ct f_0=\refl{y_0}.\]
\end{defn}

带点映射上还有另一个函子，它将 $f:X\to Y$ 映到 $\proj1 : \hfib f {y_0} \to X$。
当 $f$ 是带点映射时，我们总是将 $\hfib f {y_0}$ 视为以 $(x_0,f_0)$ 为基点的带点类型；在这种情况下，$\proj1$ 也是一个带点映射，其证据为 $(\proj1)_0 \defeq \refl{x_0}$。
因此，这个操作可以迭代进行。

\begin{defn}
  一个带点映射 $f:X\to Y$ 的 \define{纤维列（fiber sequence）}
  \indexdef{fiber sequence}%
  \indexdef{sequence!fiber}%
  是由带点类型与带点映射组成的无穷序列
  \[\xymatrix{\dots \ar[r]^{f^{(n+1)}} & X^{(n+1)} \ar[r]^{f^{(n)}} & X^{(n)} \ar^-{f^{(n-1)}}[r] & \dots \ar[r] & X^{(2)} \ar^-{f^{(1)}}[r] & X^{(1)} \ar[r]^{f^{(0)}} & X^{(0)}}\]
  递归地由
  \[ X^{(0)} \defeq Y\qquad
    X^{(1)} \defeq X\qquad
    f^{(0)} \defeq f\qquad
    \]
  与
  \begin{alignat*}{2}
    X^{(n+1)} &\defeq \hfib {f^{(n-1)}}{x^{(n-1)}_0}\\
    f^{(n)} &\defeq \proj1 &: X^{(n+1)} \to X^{(n)}.
  \end{alignat*}
  定义；
  其中 $x^{(n)}_0$ 表示 $X^{(n)}$ 的基点，按上述方式递归地选取。
\end{defn}

因此，这个纤维列中任意相邻的一对映射都具有如下形式：
\[ \xymatrix{ X^{(n+1)} \jdeq \hfib{f^{(n-1)}}{x^{(n-1)}_0} \ar[rr]^-{f^{(n)}\jdeq \proj1} && X^{(n)} \ar[r]^{f^{(n-1)}} & X^{(n-1)}. } \]
特别地，我们有 $f^{(n-1)} \circ f^{(n)} = 0$。
现在我们注意到，这个序列中出现的类型分别是底空间 $Y$、全空间\index{total!space} $X$ 以及纤维 $F\defeq \hfib f {y_0}$ 的迭代环路空间\index{loop space!iterated}，映射也类似。

\begin{lem}\label{thm:fiber-of-the-fiber}
  设 $f:X\to Y$ 是带点空间之间的一个带点映射。那么：
  \begin{enumerate}
  \item $f^{(1)}\defeq \proj1 : \hfib f {y_0} \to X$ 的纤维等价于 $\Omega Y$。\label{item:fibseq1}
  \item 类似地，$f^{(2)} : \Omega Y \to \hfib f {y_0}$ 的纤维等价于 $\Omega X$。\label{item:fibseq2}
  \item 在这些等价之下，带点映射 $f^{(3)} : \Omega X\to \Omega Y$ 等同于带点映射 $\Omega f\circ\rev{(\blank)}$。\label{item:fibseq3}
  \end{enumerate}
\end{lem}

\begin{proof}
  对于\ref{item:fibseq1}，我们有
  \begin{align*}
    \hfib{f^{(1)}}{x_0}
    &\defeq \sm{z:\hfib{f}{y_0}} (\proj{1}(z) = x_0)\\
    &\eqvsym \sm{x:X}{p:f(x)=y_0} (x = x_0) &\text{(by \cref{ex:sigma-assoc})}\\
    &\eqvsym (f(x_0) = y_0) &\text{(as $\tsm{x:X} (x=x_0)$ is contractible)}\\
    &\eqvsym (y_0 = y_0) &\text{(by $(f_0 \ct \blank)$)}\\
    &\jdeq \Omega Y.
  \end{align*}。
  追踪计算过程可知，这个等价将 $((x,p),q)$ 映到 $\opp{f_0} \ct \ap{f}{\opp q} \ct p$，而其逆将 $r:y_0=y_0$ 映到 $((x_0, f_0 \ct r), \refl{x_0})$。
  特别地，$\hfib{f^{(1)}}{x_0}$ 的基点 $((x_0,f_0),\refl{x_0})$ 被映到 $\opp{f_0} \ct \ap{f}{\opp {\refl{x_0}}} \ct f_0$，后者等于 $\refl{y_0}$。
  因此这个等价是一个带点映射（参见 \cref{ex:pointed-equivalences}）。
  而且，在这个等价之下，$f^{(2)}$ 被等同于 $\lam{r} (x_0, f_0 \ct r): \Omega Y \to \hfib f {y_0}$。

  \cref{item:fibseq2} 通过将\ref{item:fibseq1} 应用于 $f^{(1)}$（代替 $f$）立即得到。
  % The resulting equivalence sends $(((x,p),p'),q') : \hfib{f^{(2)}}{(x_0,f_0)}$ to $\ap{f}{\opp {q'}} \ct p'$, while its inverse sends $s:x_0=x_0$ to $(((x_0,f_0), s), \refl{(x_0,f_0)})$.
  由于 $(f^{(1)})_0 \defeq \refl{x_0}$，在这个等价下 $f^{(3)}$ 被等同于映射 $\Omega X \to \hfib {f^{(1)}}{x_0}$，后者由 $s \mapsto ((x_0,f_0),s)$ 定义。
  因此，当我们与前述等价 $\hfib {f^{(1)}}{x_0} \eqvsym \Omega Y$ 复合时，可以看到 $s$ 被映到 $\opp{f_0} \ct \ap{f}{\opp s} \ct f_0$，根据定义这就是 $(\Omega f)(\opp s)$。我们略去以下证明：这是带点映射之间的相等，而不仅仅是函数之间的相等。
\end{proof}

于是，$f:X\to Y$ 的纤维列可以图示如下：
\[\xymatrix@C=1.5pc{
  \dots \ar[r] &
  \Omega^2 X \ar[r]^-{\Omega^2 f} &
  \Omega^2 Y \ar[r]^-{-\Omega \partial} &
  \Omega F \ar[r]^-{-\Omega i} &
  \Omega X \ar[r]^-{-\Omega f} &
  \Omega Y \ar[r]^-{\partial} &
  F \ar[r]^-{i} &
  X \ar[r]^{f} & Y.}\]
其中减号表示与道路逆 $\rev{(\blank)}$ 的复合。
注意，由 \cref{ex:ap-path-inversion} 可得
\[ \Omega\left(\Omega f\circ \opp{(\blank)}\right) \circ \opp{(\blank)}
= \Omega^2 f \circ \opp{(\blank)} \circ \opp{(\blank)}
= \Omega^2 f.
\]
因此，当 $k$ 是奇数时，$k$ 重环路映射上就会出现减号。

从这个纤维列我们将推导出一个\emph{带点集的正合列}。
\index{image}%
设 $A$ 和 $B$ 是集合，$f:A\to B$ 是一个函数；回顾 \emph{像} $\im(f)$ 的定义（参见 \cref{defn:modal-image}），它可以视为 $B$ 的一个子集：
\[\im(f) \defeq \setof{b:B | \exis{a:A} f(a)=b}. \]
如果 $A$ 和 $B$ 还分别带有基点 $a_0$ 和 $b_0$，并且 $f$ 是一个带点映射，则定义 $f$ 的\define{核（kernel）}
\indexdef{kernel}%
\indexdef{pointed!map!kernel of}%
为 $A$ 的如下子集：
\symlabel{kernel}
\[\ker(f) \defeq \setof{x:A | f(x) = b_0}. \]
当然，这其实就是 $f$ 在基点 $b_0$ 上的纤维；由于 $B$ 是集合，所以它是 $A$ 的子集。

注意，任何群都是一个带点集，其单位元即为基点；任何群同态都是一个带点映射。
在这种情形下，核与像的定义与通常群论中的概念一致。

\begin{defn}
  一个\define{带点集的正合列（exact sequence of pointed sets）}
  \indexdef{exact sequence}%
  \indexdef{sequence!exact}%
  是一个（可能有界的）带点集与带点映射的序列：
  \[\xymatrix{\dots \ar[r] & A^{(n+1)} \ar[r]^-{f^{(n)}} & A^{(n)} \ar[r]^{f^{(n-1)}} & A^{(n-1)} \ar[r] &
    \dots}\]
  使得对每个 $n$，$f^{(n)}$ 的像作为 $A^{(n)}$ 的子集等于 $f^{(n-1)}$ 的核。
  换言之，对所有 $a:A^{(n)}$，有
  \[ (f^{(n-1)}(a) = a^{(n-1)}_0) \iff \exis{b:A^{(n+1)}} (f^{(n)}(b)=a). \]
  其中 $a^{(n)}_0$ 表示 $A^{(n)}$ 的基点。
\end{defn}

通常，正合列中的大多数或全部带点集都是群，且常为交换群。
当我们提到\define{群的正合列（exact sequence of groups）}时，还假定其中的映射是群同态，而不仅仅是带点映射。

\index{exact sequence}


\begin{thm}\label{thm:les}
  设 $f:X \to Y$ 是带点空间之间的一个带点映射，其纤维为 $F\defeq \hfib f {y_0}$。
  则我们得到如下长正合列，其中除最后三项外都是群；除最后六项外，都是交换群。

  \[
  \xymatrix@R=1.2pc@C=3pc{
    \vdots & \vdots & \vdots \ar[lld] \\
    \pi_k(F) \ar[r] & \pi_k(X) \ar[r] & \pi_k(Y) \ar[lld] \\
    \vdots & \vdots & \vdots \ar[lld] \\
    \pi_2(F) \ar[r] & \pi_2(X) \ar[r] & \pi_2(Y) \ar[lld] \\
    \pi_1(F) \ar[r] & \pi_1(X) \ar[r] & \pi_1(Y) \ar[lld]\\
    \pi_0(F) \ar[r] & \pi_0(X) \ar[r] & \pi_0(Y)}
  \]
\end{thm}

% In US trade we might want a page break here, and extra stretch, otherwise the
% whole page looks ugly.
\vspace*{0pt plus 10ex}
\goodbreak

\begin{proof}
  我们首先证明，一个纤维列的 $0$-截断是一个带点集的正合列。
  因此，对于纤维列中任意相邻的一对映射：
  %
  \[\xymatrix{\hfib{f}{z_0} \ar^-g[r] & W \ar^-f[r] & Z}\]
  %
  其中 $g\defeq \proj1$，我们需要证明列
  %
  \[\xymatrix{\trunc{0}{\hfib{f}{z_0}} \ar^-{\trunc0g}[r] & \trunc0{W}
    \ar^-{\trunc0{f}}[r] & \trunc0{Z}}\]
  %
  是正合的，即证明 $\im(\trunc0g)\subseteq\ker(\trunc0f)$ 与 $\ker(\trunc0f)\subseteq\im(\trunc0g)$。

  第一个包含关系等价于 $\trunc0f\circ\trunc0g=0$，这由 $\truncf0$ 的函子性以及 $f\circ g=0$ 可得。
  对于第二个包含关系，我们假设 $w':\trunc0W$ 且 $p':\trunc0f(w')=\tproj0{z_0}$，并证明仅存在 ${t:\hfib{f}{z_0}}$ 使得 $\tproj0{g(t)}=w'$。
  由于我们的目标是一个命题，可以假定 $w'$ 具有 $\tproj0w$ 的形式，其中 $w:W$。
  现在，根据 \cref{thm:path-truncation}，由 $p' : \tproj0{f(w)}=\tproj0{z_0}$ 可得 $p'': \trunc{-1}{f(w)=z_0}$，因此通过进一步的截断归纳，我们可以假设存在某个 $p:f(w)=z_0$。
  但这样我们就得到了 $\tproj0{(w,p)}:\trunc{0}{\hfib{f}{z_0}}$，其在 $\trunc0g$ 下的像正是 $\tproj0w\jdeq w'$，如所需。

  因此，将 $\truncf0$ 应用于 $f$ 的纤维列，我们得到了一个长正合列，其中按所需顺序包含了带点集 $\pi_k(F)$、$\pi_k(X)$ 和 $\pi_k(Y)$。
  当然，对于 $k\ge1$，$\pi_k$ 是一个群（因为它是环路空间的 $0$-截断），并且对于 $k\ge 2$，由 Eckmann--Hilton 论证可知它是交换群
  \index{Eckmann--Hilton argument} (\cref{thm:EckmannHilton})。
  此外，\cref{thm:fiber-of-the-fiber} 允许我们将这个正合列中的映射 $\pi_k(F) \to \pi_k(X)$ 和 $\pi_k(X) \to \pi_k(Y)$ 分别识别为 $(-1)^k \pi_k(i)$ 和 $(-1)^k \pi_k(f)$。

  更一般地，这个长正合列中除最后三个映射外，每个映射都具有 $\trunc0{\Omega h}$ 或 $\trunc0{-\Omega h}$ 的形式，其中 $h$ 是某个映射。
  在前一种情形下，它是群同态；而在后一种情形下，若所涉及的群是交换群，它也是同态；否则它是一个``反同态''。
  然而，当我们把一个群同态替换为它的负映射时，其核与像保持不变，因此任何涉及该映射的列的正合性也保持不变。
  于是，我们可以修改我们的长正合列，得到一个直接包含 $\pi_k(i)$ 和 $\pi_k(f)$ 的列，并且其中所有映射（除最后三个外）都是群同态。
\end{proof}

交换群的正合列的通常性质\index{group!abelian!exact sequence of}可以按通常方式证明。
特别地，我们有：


\begin{lem}\label{thm:ses}
  给定一个交换群的正合列：
  %
  \[\xymatrix{K \ar[r]& G \ar^f[r] & H \ar[r] & Q.}\]
  %
  \begin{enumerate}
  \item 若 $K=0$，则 $f$ 是单的。\label{item:sesinj}
  \item 若 $Q=0$，则 $f$ 是满的。\label{item:sessurj}
  \item 若 $K=Q=0$，则 $f$ 是同构。\label{item:sesiso}
  \end{enumerate}
\end{lem}

\begin{proof}
  由于 $f$ 的核是 $K\to G$ 的像，若 $K=0$，则 $f$ 的核是 $\{0\}$；
  因此 $f$ 是单的，因为它是群同态。
  类似地，由于 $f$ 的像是 $H\to Q$ 的核，若 $Q=0$，则 $f$ 的像是整个 $H$，所以 $f$ 是满的。
  最后，\ref{item:sesiso} 由 \ref{item:sesinj} 与 \ref{item:sessurj} 根据 \cref{thm:mono-surj-equiv} 得出。
\end{proof}

作为一个直接应用，我们现在可以具体说明，一个映射的 $n$-连通性在何种意义上比它在 $n$-截断上诱导等价更强。

\begin{cor}\label{thm:conn-pik}
  设 $f:A\to B$ 是 $n$-连通的，$a:A$，并定义 $b\defeq f(a)$。那么：
  \begin{enumerate}
  \item 若 $k\le n$，则 $\pi_k(f):\pi_k(A,a) \to \pi_k(B,b)$ 是同构。\label{item:conn-pik1}
  \item 若 $k=n+1$，则 $\pi_k(f):\pi_k(A,a) \to \pi_k(B,b)$ 是满的。\label{item:conn-pik2}
  \end{enumerate}
\end{cor}

\begin{proof}
  当 $k=0$ 时，第~\ref{item:conn-pik1}部分可由~\cref{lem:connected-map-equiv-truncation}得出，只需注意到 $\pi_0(f)\equiv\trunc 0f$。
  当 $k=0$ 时，第~\ref{item:conn-pik2}部分可由~\cref{ex:is-conn-trunc-functor}得出，注意到一个函数是满的当且仅当它是 $(-1)$-连通的，依据~\cref{thm:minusoneconn-surjective}。
  当 $k>0$ 时，作为长正合列的一部分，我们得到一个正合列
  \[\xymatrix{\pi_k(\hfib f b) \ar[r]& \pi_k(A,a) \ar^f[r] & \pi_k(B,b) \ar[r] & \pi_{k-1}(\hfib f b).}\]
  现在，由于 $f$ 是 $n$-连通的，$\trunc n{\hfib f b}$ 是可缩的。
  因此，若 $k\le n$，则~$\pi_k(\hfib f b) = \trunc0{\Omega^k(\hfib f b)} = \Omega^k(\trunc k{\hfib f b})$也是可缩的。
  这样一来，由~\cref{thm:ses}\ref{item:sesiso}可知，当 $k\le n$ 时 $\pi_k(f)$ 是同构；而当 $k=n+1$ 时，由~\cref{thm:ses}\ref{item:sessurj}可知它是满的。
\end{proof}

在~\cref{sec:whitehead}中我们将会看到，\cref{thm:conn-pik}的逆命题也成立。

\index{fiber sequence|)}%
\index{sequence!fiber|)}%

\section{Hopf 纤维化}
\label{sec:hopf}

在本节中，我们将定义\define{Hopf 纤维化（Hopf fibration）}。
\indexdef{Hopf!fibration}%

\begin{thm}[Hopf 纤维化]\label{thm:hopf-fibration}
  存在 $\Sn ^2$ 上的一个纤维化 $H$，它在基点处的纤维是 $\Sn ^1$，全空间是 $\Sn ^3$。
\end{thm}

Hopf 纤维化将帮助我们计算球面的若干同伦群。
事实上，它给出了下面的同伦群长正合列
\index{homotopy!group!of sphere}
\index{sequence!exact}
（参见
\cref{sec:long-exact-sequence-homotopy-groups}）：
%
\[
\xymatrix@R=1.2pc{
  \pi_k(\Sn^1) \ar[r] & \pi_k(\Sn^3) \ar[r] & \pi_k(\Sn^2) \ar[lld] \\
  \vdots & \vdots & \vdots \ar[lld] \\
  \pi_2(\Sn^1) \ar[r] & \pi_2(\Sn^3) \ar[r] & \pi_2(\Sn^2) \ar[lld] \\
  \pi_1(\Sn^1) \ar[r] & \pi_1(\Sn^3) \ar[r] & \pi_1(\Sn^2)}
\]
%
我们已经计算了所有的 $\pi_n(\Sn^1)$，以及当 $k<n$ 时的 $\pi_k(\Sn^n)$，因此上式化为：
%
\[
\xymatrix@R=1.2pc{
  0 \ar[r] & \pi_k(\Sn^3) \ar[r] & \pi_k(\Sn^2) \ar[lld] \\
  \vdots & \vdots & \vdots \ar[lld] \\
  0 \ar[r] & \pi_3(\Sn^3) \ar[r] & \pi_3(\Sn^2) \ar[lld] \\
  0 \ar[r] & 0 \ar[r] & \pi_2(\Sn^2) \ar[lld] \\
  \Z \ar[r] & 0 \ar[r] & 0}
\]
%
特别地，我们得到如下结果：

\begin{cor} \label{cor:pis2-hopf}
  对于每个 $k\ge3$，有 $\eqv{\pi_2(\Sn^2)}{\Z}$ 和 $\eqv{\pi_k(\Sn^3)}{\pi_k(\Sn^2)}$（其中该映射由 Hopf 纤维化诱导，视为从全空间 $\Sn^3$ 到基空间 $\Sn^2$ 的映射）。
\end{cor}

事实上，我们可以更进一步：Hopf 纤维化的纤维列将表明 $\Omega^3(\Sn^3)$ 是从 $\Omega^3(\Sn^2)$ 到 $\Omega^2(\Sn^1)$ 的某个映射的纤维。
由于 $\Omega^2(\Sn^1)$ 是可缩的，我们有 $\eqv{\Omega^3(\Sn^3)}{\Omega^3(\Sn^2)}$。
在经典同伦论中，这个事实会是 \cref{cor:pis2-hopf} 和 Whitehead 定理的推论，但 Whitehead 定理在同伦类型论中不一定成立（参见 \cref{sec:whitehead}）。
不过，我们在这里不会使用更精确的版本。

\subsection{推出上的纤维化}
\label{sec:fib-over-pushout}

我们首先给出一个引理，说明如何在推出上构造纤维化。
\index{pushout}%

\begin{lem}\label{lem:fibration-over-pushout}
  设 $\Ddiag=(Y\xleftarrow{j}X\xrightarrow{k}Z)$ 是一个伸展\index{span}，并假设我们有：
  \begin{itemize}
  \item 两个纤维化 $E_Y:Y\to\type$ 和 $E_Z:Z\to\type$。
  \item $E_Y\circ j:X\to\type$ 与 $E_Z\circ k:X\to\type$ 之间的一个等价 $e_X$，即
    \[e_X:\prd{x:X}\eqv{E_Y(j(x))}{E_Z(k(x))}.\]
  \end{itemize}

  那么我们可以构造一个纤维化 $E:Y\sqcup^XZ\to\type$，使得：
  \begin{itemize}
  \item 对于所有 $y:Y$，$E(\inl(y))\judgeq E_Y(y)$。
  \item 对于所有 $z:Z$，$E(\inr(z))\judgeq E_Z(z)$。
  \item 对于所有 $x:X$，$\map E{\glue(x)}=\ua(e_X(x))$（注意等式两边都是 $\type$ 中从 $E_Y(j(x))$ 到 $E_Z(k(x))$ 的道路）。
  \end{itemize}
  此外，该纤维化的全空间构成如下推出方块：
  \[\xymatrix{ \sm{x:X}E_Y(j(x)) \ar[r]_\sim^{\idfunc\times e_X}
    \ar[d]_{j\times\idfunc} &
    \sm{x:X}E_Z(k(x)) \ar[r]^-{k\times\idfunc}
    & \sm{z:Z}E_Z(z) \ar[d]^{\inr} \\
    \sm{y:Y}E_Y(y) \ar[rr]_{\inl} & & \sm{t:Y\sqcup^XZ}E(t) }\]
\end{lem}

\begin{proof}
  我们利用推出 $Y\sqcup^XZ$ 的递归原理来定义 $E$。
  为此，我们需要指定 $E$ 在形如 $\inl(y)$、$\inr(z)$ 的元素上的值，以及 $E$ 在道路 $\glue(x)$ 上的作用，因此可以直接选取以下值：
  \begin{align*}
    E(\inl(y)) & \defeq E_Y(y),\\
    E(\inr(z)) & \defeq E_Z(z),\\
    \map E{\glue(x)} & \defid \ua(e_X(x)).
  \end{align*}
  %
  为证明该纤维化的全空间是一个推出，我们应用展平引理（\cref{thm:flattening}）并取以下值：\index{flattening lemma}
  \begin{itemize}
  \item $A\defeq Y+Z$，$B\defeq X$，且 $f,g:B\to A$ 定义为
    $f(x)\defeq\inl(j(x))$，$g(x)\defeq\inr(k(x))$，
  \item 类型族 $C:A\to\type$ 定义为
    \begin{equation*}
      C(\inl(y)) \defeq E_Y(y)
      \qquad\text{and}\qquad
      C(\inr(z)) \defeq E_Z(z),
    \end{equation*}
  \item 等价族 $D:\prd{b:B}C(f(b))\eqvsym C(g(b))$ 定义为 $e_X$。
  \end{itemize}
  %
  展平引理中的基础高阶归纳类型 $W$ 等价于推出 $Y\sqcup^XZ$，而类型族 $P:Y\sqcup^XZ\to\type$ 等价于上面定义的 $E$。

  因此展平引理告诉我们 $\sm{t:Y\sqcup^XZ}E(t)$ 等价于具有以下生成元的高阶归纳类型 ${E^{\mathrm{tot}}}'$：
  %
  \begin{itemize}
  \item 一个函数 $\mathsf{z}:\sm{a:Y+Z}C(a)\to {E^{\mathrm{tot}}}'$，
  \item 对每个 $x:X$ 和 $t:E_Y(j(x))$，一条道路
    \narrowequation{\mathsf{z}(\inl(j(x)),t)=\mathsf{z}(\inr(k(x)),e_X(t))。}
  \end{itemize}
  %
  再次使用展平引理或直接计算，容易看出 $\eqv{\sm{a:Y+Z}C(a)}{\sm{y:Y}E_Y(y)+\sm{z:Z}E_Z(z)}$，因此
  ${E^{\mathrm{tot}}}'$ 等价于具有以下生成元的高阶归纳类型 $E^{\mathrm{tot}}$：
  %
  \begin{itemize}
  \item 一个函数 $\inl:\sm{y:Y}E_Y(y)\to E^{\mathrm{tot}}$，
  \item 一个函数 $\inr:\sm{z:Z}E_Z(z)\to E^{\mathrm{tot}}$，
  \item 对每个 $(x,t):\sm{x:X}E_Y(j(x))$，一条道路
    \narrowequation{\glue(x,t):\inl(j(x),t) = \inr(k(x),e_X(t))。}
  \end{itemize}
  %
  因此，$E$ 的全空间是 $E_Y$ 和 $E_Z$ 的全空间的推出，正如所需。
\end{proof}

\subsection{Hopf 构造}

\begin{defn}
  一个 \define{H-空间（H-space）}
  \indexdef{H-space}%
  由以下结构组成：
  \begin{itemize}
  \item 一个类型 $A$，
  \item 一个基点 $e:A$，
  \item 一个二元运算 $\mu:A\times A\to A$，以及
  \item 对于每个 $a:A$，有相等 $\mu(e,a)=a$ 和 $\mu(a,e)=a$。
  \end{itemize}
\end{defn}

\begin{lem}
  设 $A$ 是一个连通的 H-空间。那么对于每个 $a:A$，映射 $\mu(a,\blank):A\to
  A$ 和 $\mu(\blank,a):A\to A$ 都是等价。
\end{lem}

\begin{proof}
  我们来证明对于每个 $a:A$，映射 $\mu(a,\blank)$ 是一个等价。另一个结论同理可证。
  %
  ``$\mu(a,\blank)$ 是等价''这一断言对应于一个类型族
  $P:A\to\prop$，而证明它则对应于寻找该类型族的一个截面。

  类型 $\prop$ 是一个集合（\cref{thm:hleveln-of-hlevelSn}），因此我们可以
  通过 $P'(\tproj0a)\defeq
  P(a)$ 定义一个新的类型族 $P':\trunc0A\to\prop$。但根据假设，$A$ 是连通的，因此 $\trunc0A$
  是可缩的。这蕴含着，要找到 $P'$ 的一个截面，只需在 $P'$ 于 $\tproj0e$ 处的纤维中找到一个点即可。而我们有
  $P'(\tproj0e)=P(e)$，该类型是有实例的，因为根据 H-空间的定义，$\mu(e,\blank)$ 等于恒等映射，因而是一个等价。

  我们已经证明了对于每个 $x:\trunc0A$，命题 $P'(x)$ 为真，
  因此特别地，对于每个 $a:A$，命题 $P(a)$ 也为真，因为
  $P(a)$ 就是 $P'(\tproj0a)$。
\end{proof}

\begin{defn}
  设 $A$ 是一个连通的 H-空间。我们利用
  \cref{lem:fibration-over-pushout} 在 $\susp A$ 上构造一个纤维化。

  鉴于 $\susp A$ 是推出 $\unit\sqcup^A\unit$，我们可以通过指定以下内容来定义 $\susp A$ 上的一个纤维化：
  \begin{itemize}
  \item $\unit$ 上的两个纤维化（即两个类型 $F_1$ 和 $F_2$），以及
  \item 一族 $e:A\to(\eqv{F_1}{F_2})$ $F_1$ 与 $F_2$ 之间的等价，$A$ 中每个元素各对应一个。
  \end{itemize}
  %
  我们取 $F_1$ 和 $F_2$ 都为 $A$，并对 $a:A$ 取等价 $\mu(a,\blank)$ 作为 $e(a)$。
\end{defn}

根据 \cref{lem:fibration-over-pushout}，我们得到如下图表：
%
\[\xymatrix{A \ar@{->>}[d] & A \times A \ar[l]_-{\proj2} \ar@{->>}_{\proj1}[d]
  \ar[r]^-{\mu} & A \ar@{->>}[d] \\
  1 & A \ar[r] \ar[l] & 1}\]
%
我们刚才构造的纤维化是 $\susp A$ 上的一个纤维化，其全空间是上方那一行的推出。

此外，利用 $f(x,y)\defeq(\mu(x,y),y)$ 我们得到如下图表：
%
\[\xymatrix{A \ar_{\idfunc}[d] & A \times A \ar[l]_-{\proj2} \ar^f[d]
  \ar[r]^-{\mu} & A \ar^{\idfunc}[d] \\
  A & A\times A \ar^-{\proj2}[l] \ar_-{\proj1}[r] & A}\]
%
该图表交换，并且三个垂直映射都是等价，其中 $f$ 的逆是函数 $g$，定义为
\[g(u,v)\defeq(\opp{\mu(\blank,v)}(u),v).\]
%
这表明两行是等价的（从而相等的）伸展，因此我们所构造纤维化的全空间等价于下方那一行的推出。
而根据定义，这个推出就是 $A$ 与自身的\emph{联结}（见 \cref{sec:colimits}）。
%
于是我们证明了：

\begin{lem}\label{lem:hopf-construction}
  给定一个连通的 H-空间 $A$，存在一个纤维化，称为
  \define{Hopf 构造（Hopf construction）}，
  \indexdef{Hopf!construction}%
  其底空间为 $\susp A$，纤维为 $A$，全空间为 $A*A$。
\end{lem}

\subsection{Hopf 纤维化}

\index{Hopf fibration|(}
\indexsee{fibration!Hopf}{Hopf fibration}
我们将首先在圆 $\Sn^1$ 上构造一个 H-空间结构，从而由
\cref{lem:hopf-construction} 我们将得到一个底空间为 $\Sn^2$，纤维为
$\Sn^1$ 且全空间为 $\Sn^1*\Sn^1$ 的纤维化。然后我们将证明这个联结等价于 $\Sn^3$。

\begin{lem}\label{lem:hspace-S1}
  在圆 $\Sn^1$ 上存在一个 H-空间结构。
\end{lem}

\begin{proof}
  对于 H-空间结构的基点，我们选择 $\base$。
  %
  现在我们需要定义乘法运算
  $\mu:\Sn^1\times\Sn^1\to\Sn^1$。
  我们将通过对 $\Sn^1$ 作递归来定义 $\mu$ 的柯里化形式 $\widetilde\mu:\Sn^1\to(\Sn^1\to\Sn^1)$：
  \begin{equation*}
    \widetilde\mu(\base) \defeq\idfunc[\Sn^1],
    \qquad\text{and}\qquad
    \ap{\widetilde\mu}{\lloop} \defid\funext(h).
  \end{equation*}
  其中 $h:\prd{x:\Sn^1}(x=x)$ 是 \cref{thm:S1-autohtpy} 中定义的函数，
  它具有性质 $h(\base) \defeq\lloop$。

  现在只需证明对每个 $x:\Sn^1$，有 $\mu(x,\base)=\mu(\base,x)=x$。
  %
  根据定义，对于 $x:\Sn^1$，我们有
  $\mu(\base,x)=\widetilde\mu(\base)(x)=\idfunc[\Sn^1](x)=x$。为了证明等式
  $\mu(x,\base)=x$，我们对 $x:\Sn^1$ 进行归纳：
  \begin{itemize}
  \item 如果 $x$ 是 $\base$，那么根据定义 $\mu(\base,\base)=\base$，因此我们有 $\refl\base:\mu(\base,\base)=\base$。
  \item 当 $x$ 沿 $\lloop$ 变化时，我们需要证明
    \[\refl\base\ct\apfunc{\lam{x}x}({\lloop})=
    \apfunc{\lam{x}\mu(x,\base)}({\lloop})\ct\refl\base.\]
    左边等于 $\lloop$，而对于右边我们有：
    \begin{align*}
      \apfunc{\lam{x}\mu(x,\base)}({\lloop})\ct\refl\base &=
      \apfunc{\lam{x}(\widetilde\mu(x))(\base)}({\lloop})\\
      &=\happly(\apfunc{\lam{x}(\widetilde\mu(x))}({\lloop}),\base)\\
      &=\happly(\funext(h),\base)\\
      &=h(\base)\\
      &=\lloop. \qedhere
    \end{align*}
  \end{itemize}
\end{proof}

现在回顾 \cref{sec:colimits}，类型 $A$ 和 $B$ 的\emph{联结} $A*B$ 是图表
\index{join!of types}%
\[A \xleftarrow{\proj1}A\times B \xrightarrow{\proj2} B. \] 的推出。

\begin{lem}
  \index{associativity!of join}%
  联结运算是结合的：若 $A$、$B$ 和 $C$ 是三个类型，
  则我们有等价 $\eqv{(A*B)*C}{A*(B*C)}$。
\end{lem}

\begin{proof}
  我们通过归纳来定义一个映射 $f:(A*B)*C\to A*(B*C)$。首先需要定义
  $f\circ\inl:A*B\to A*(B*C)$，这也将通过归纳来完成，然后定义
  $f\circ\inr:C\to A*(B*C)$，接着定义 $\apfunc{f}\circ\glue:\prd{t:(A*B)\times
    C}f(\inl(\fst(t)))=f(\inr(\snd(t)))$，后者需要对 $t$ 的第一个分量进行归纳：
  \begin{align*}
    (f\circ\inl)(\inl(a)) &\defeq \inl(a), \\
    (f\circ\inl)(\inr(b)) &\defeq \inr(\inl(b)), \\
    \apfunc{f\circ\inl}(\glue(a,b)) &\defid \glue(a,\inl(b)), \\
    f(\inr(c)) &\defeq \inr(\inr(c)),\\
    \apfunc{f}(\glue(\inl(a),c)) &\defid\glue(a,\inr(c)),\\
    \apfunc{f}(\glue(\inr(b),c)) &\defid\apfunc{\inr}(\glue(b,c)),\\
    \apdfunc{\lam{x}\apfunc{f}(\glue(x,c))}(\glue(a,b)) &\defid
    ``\apdfunc{\lam{x}\glue(a,x)}(\glue(b,c))''.
  \end{align*}
  %
  对于最后一个等式，注意其右侧的类型是
  \[\transfib{\lam{x}\inl(a)=\inr(x)}{\glue(b,c)}{\glue(a,\inl(b))}=
  \glue(a,\inr(c))\]
  而它本应具有类型
  %
  \begin{narrowmultline*}
    \transfib{\lam{x}f(\inl(x))=f(\inr(c))}{\glue(a,b)}{\apfunc{f}(\glue(\inl(a),c))}
    = \narrowbreak
    \apfunc{f}(\glue(\inr(b),c)).
  \end{narrowmultline*}
  %
  但根据定义中前面的几条，这两个类型都等价于以下类型：
  \[\glue(a,\inr(c))=\glue(a,\inl(b))\ct\apfunc\inr(\glue(b,c)),\]
  因此，我们可以通过一个等价进行强制转换，从而得到所需的元素。
  %
  类似地，我们可以定义一个映射 $g:A*(B*C)\to(A*B)*C$，而验证 $f$ 和
  $g$ 互为逆映射，则是一段冗长繁琐但在本质上直截了当的计算。
\end{proof}

一个更具概念性的证明概要如下。

\begin{proof}
  考虑如下图表，其中的映射都是自然的投影：
  \[\xymatrix{
    A & A\times C \ar[l] \ar[r] & A\times C\\
    A\times B \ar[u] \ar[d] & A\times B\times C \ar[l]\ar[u]\ar[r]\ar[d] &
    A\times C \ar[u] \ar[d] \\
    B & B\times C \ar[l] \ar[r] & C}\]
  %
  对各列取余极限，得到如下图表，其余极限是 $(A*B)*C$：
  \[\xymatrix{A*B & (A*B)\times C \ar[l]\ar[r] & C}\]
  %
  另一方面，对各行取余极限，得到一个图表，其余极限是 $A*(B*C)$。

  因此，利用一个关于余极限的类 Fubini 定理（我们尚未证明），便可得到等价 $\eqv{(A*B)*C}{A*(B*C)}$。不过，这个关于余极限的 Fubini 定理
  \index{Fubini theorem for colimits}
  其证明仍需那段冗长繁琐的计算。
\end{proof}

\begin{lem}
  对于任意类型 $A$，存在等价 $\eqv{\susp A}{\bool*A}$。
\end{lem}

\begin{proof}
  定义来回两个方向的映射很容易，证明它们互为逆映射也很直接。详情留给读者作为练习。
\end{proof}

现在我们可以构造 Hopf 纤维化：

\begin{thm}
  存在一个以 $\Sn^2$ 为底、$\Sn^1$ 为纤维、$\Sn^3$ 为全空间的纤维化。
  \index{total!space}%
\end{thm}

\begin{proof}
  我们先前证明了 $\Sn^1$ 具有 H-空间结构（参见 \cref{lem:hspace-S1}），
  因此根据 \cref{lem:hopf-construction}，存在一个以 $\Sn^2$ 为底、$\Sn^1$ 为纤维、$\Sn^1*\Sn^1$ 为全空间的纤维化。
  但由前两个结果以及 \cref{thm:suspbool}，我们有：
  \begin{equation*}
    \Sn^1*\Sn^1 = (\susp\bool)*\Sn^1
    =(\bool*\bool)*\Sn^1
    =\bool*(\bool*\Sn^1)
    =\susp(\susp\Sn^1)
    =\Sn^3. \qedhere
  \end{equation*}
\end{proof}


\index{Hopf fibration|)}

\section{Freudenthal 纬悬定理}
\label{sec:freudenthal}

\index{Freudenthal suspension theorem|(}%
\index{theorem!Freudenthal suspension|(}%

在证明 Freudenthal 纬悬定理之前，我们需要一些关于连通性的辅助引理。
在\cref{cha:hlevels}中，我们针对固定的 $n$ 证明了关于 $n$-连通映射与 $n$-类型的一系列事实；现在我们关心的是当 $n$ 变化时会发生什么。
例如，在\cref{prop:nconnected_tested_by_lv_n_dependent types}中我们证明了 $n$-连通映射可由一个相对于 $n$-类型族的``归纳原理''来刻画。
如果我们想沿着一个 $n$-连通映射，向一个 $k>n$ 的 $k$-类型族作归纳，我们并不能立即知道存在一个满足该归纳原理的函数，但下面的引理表明，至少我们的无知是可以被量化的。

\begin{lem}\label{thm:conn-trunc-variable-ind}
  若 $f:A\to B$ 是 $n$-连通的，并且 $P:B\to \ntype{k}$ 是 $k$-类型族（其中 $k\ge n$），则诱导函数
  \[ (\blank\circ f) : \Parens{\prd{b:B} P(b)} \to \Parens{\prd{a:A} P(f(a)) } \]
  是 $(k-n-2)$-截断的。
\end{lem}

\begin{proof}
  我们对自然数 $k-n$ 进行归纳。
  当 $k=n$ 时，这就是\cref{prop:nconnected_tested_by_lv_n_dependent types}。
  %
  对于归纳步骤，假设 $f$ 是 $n$-连通的且 $P$ 是 $(k+1)$-类型族。
  为证明 $(\blank\circ f)$ 是 $(k-n-1)$-截断的，令$\ell:\prd{a:A} P(f(a))$；则我们有
  \[ \hfib{(\blank\circ f)}{\ell} \eqvsym \sm{g:\prd{b:B} P(b)} \prd{a:A} g(f(a)) = \ell(a).\]
  设 $(g,p)$ 和 $(h,q)$ 属于此类型，即 $p:g\circ f \htpy \ell$ 且 $q:h\circ f \htpy \ell$；则我们还知道
  \[ \big((g,p) = (h,q)\big) \eqvsym
  \Parens{\sm{r:g\htpy h} r\circ f = p \ct \opp{q}}.
  \]
  然而，此处右侧是映射
  \[ (\blank\circ f) : \Parens{\prd{b:B} Q(b)} \to \Parens{\prd{a:A} Q(f(a)) } \]
  的一个纤维，其中$Q(b) \defeq (g(b)=h(b))$。
  由于 $P$ 是 $(k+1)$-类型族，$Q$ 是 $k$-类型族，故归纳假设表明该纤维是 $(k-n-2)$-类型。
  因此，$\hfib{(\blank\circ f)}{\ell}$ 的所有路径空间都是 $(k-n-2)$-类型，从而它是一个 $(k-n-1)$-类型。
\end{proof}

回顾一下，如果 $\pairr{A,a_0}$ 和 $\pairr{B,b_0}$ 是带点类型，那么它们的 \define{楔积}
\index{wedge}%
$A\vee B$ 定义为由 $A\xleftarrow{a_0}
\unit\xrightarrow{b_0} B$ 给出的推出。
存在一个标准映射 $i:A\vee B \to A\times B$，由两个映射 $\lam{a} (a,b_0)$ 和 $\lam{b} (a_0,b)$ 定义；下面的引理本质上说，如果 $A$ 和 $B$ 是高度连通的，那么这个映射也是高度连通的。
不过，如果我们使用 \cref{prop:nconnected_tested_by_lv_n_dependent types} 中连通性的刻画，并将其代入楔积的（推广到类型族的）泛性质，那么无论是证明还是使用起来都更方便一些。

\begin{lem}[楔连通性引理]\label{thm:wedge-connectivity}
  假设 $\pairr{A,a_0}$ 和 $\pairr{B,b_0}$ 分别是 $n$-连通和 $m$-连通的带点类型，其中 $n,m\geq0$，并令
%
\narrowequation{P:A\to B\to \ntype{(n+m)}.}
%
那么对于任意 ${f:\prd{a:A} P(a,b_0)}$ 与 ${g:\prd{b:B} P(a_0,b)}$，其中 $p:f(a_0) = g(b_0)$，存在 $h:\prd{a:A}{b:B} P(a,b)$ 以及同伦
%
\begin{equation*}
  q:\prd{a:A} h(a,b_0)=f(a)
  \qquad\text{and}\qquad
  r:\prd{b:B} h(a_0,b)=g(b)
 \end{equation*}
%
使得 $p = \opp{q(a_0)} \ct r(b_0)$。
\end{lem}

\begin{proof}
  定义 $Q:A\to\type$ 为
  \[ Q(a) \defeq \sm{k:\prd{b:B} P(a,b)} (f(a) = k(b_0)). \]
  则我们有 $(g,p):Q(a_0)$。
  由于 $a_0:\unit\to A$ 是 $(n-1)$-连通的，若 $Q$ 是 $(n-1)$-类型族，则我们将得到 $\ell:\prd{a:A} Q(a)$ 使得 $\ell(a_0) = (g,p)$；此时可以定义 $h(a,b) \defeq \proj1(\ell(a))(b)$。
  然而，对于固定的 $a$，类型 $Q(a)$ 是映射
  \[ \Parens{\prd{b:B} P(a,b) } \to P(a,b_0) \]
  在 $f(a)$ 处的纤维，该映射通过用 $b_0:\unit\to B$ 作前复合得到。
  由于 $b_0:\unit\to B$ 是 $(m-1)$-连通的，根据 \cref{thm:conn-trunc-variable-ind}，为使这个纤维是 $(n-1)$-截断的，只需每个类型 $P(a,b)$ 都是 $(n+m)$-类型即可，而这正是我们的假设。
\end{proof}

令 $(X,x_0)$ 为一个带点类型，并回顾 \cref{sec:suspension} 中纬悬 $\susp X$ 的定义，其构造子为 $\north,\south:\susp X$ 与 $\merid:X \to (\north=\south)$。
我们将 $\susp X$ 视为以 $\north$ 为基点的带点空间，因此有 $\Omega\susp X \defeq (\id[\susp X]\north\north)$。
于是存在一个规范映射
\begin{align*}
  \sigma &: X \to \Omega\susp X\\
  \sigma(x) &\defeq \merid(x) \ct \opp{\merid(x_0)}.
\end{align*}

\begin{rmk}
  在经典代数拓扑中，人们考虑的是\emph{约化纬悬}，其中道路 $\merid(x_0)$ 被坍缩为一个点，从而将 $\north$ 和 $\south$ 等同起来。
  约化与非约化纬悬是同伦等价的，因此在我们纯粹的同伦论视角下，这一区分是不可见的——况且高阶归纳类型本来也只允许我们通过高阶道路来``等同''点，所以在同伦类型论中并无必要考虑约化纬悬。
  然而，我们所用纬悬的``非约化性''正是 $\opp{\merid(x_0)}$ 这一（可能令人意外的）项在 $\sigma$ 的定义中出现的原因。
\end{rmk}

我们现在的目标是证明以下定理。

\begin{thm}[Freudenthal 纬悬定理]\label{thm:freudenthal}
  假设 $X$ 带点且 $n$-连通，其中 $n\geq 0$。
  则映射 $\sigma:X\to \Omega\susp(X)$ 是 $2n$-连通的。
\end{thm}

\index{encode-decode method|(}%

我们将使用编码-解码方法，但以稍有不同的方式应用它。
在之前的大多数情形中，我们通过构造一个满足 $\code(a_0)\defeq B$ 的族 $\code:A\to \type$，以及一族等价 $\decode:\prd{x:A}\code(x) \eqvsym (a_0=x)$，将某个类型的环路空间 $\Omega (A,a_0)$ 刻画为与另一个类型 $B$ 等价。
% We have also generalized it to characterize truncations of loop spaces by way of a family of equivalences $\prd{x:A}\code(x) \eqvsym \trunc n{a_0=x}$.

然而在本例中，我们希望证明 $\sigma:X\to \Omega \susp X$ 是 $2n$-连通的。
我们可以使用先前方法的截断版本（例如我们将在 \cref{sec:van-kampen} 中看到的那样）来证明 $\trunc{2n}X\to \trunc{2n}{\Omega \susp X}$ 是一个等价——但这比映射为 $2n$-连通这一断言稍弱（参见 \cref{thm:conn-pik,thm:pik-conn}）。
不过请注意，在一般情形中，为了证明 $\decode(x)$ 是一个等价，我们也可以等价地证明其纤维是可缩的，并且我们仍然可以对底类型进行归纳。
我们可以将这一点推广，用以证明到环路空间的映射的连通性，即证明其纤维的\emph{截断}是可缩的。
此外，与其分别构造 $\code$ 和 $\decode$，我们可以直接构造一族\emph{纤维之截断的代码（codes for the truncations of the fibers）}。

\begin{defn}\label{thm:freudcode}
  若 $X$ 带点且 $n$-连通（$n\geq 0$），则存在一族
  \begin{equation}
    \code:\prd{y:\susp X} (\north=y) \to \type\label{eq:freudcode}
  \end{equation}
  使得
  \begin{align}
    \code(\north,p) &\defeq \trunc{2n}{\hfib{\sigma}{p}}
    \jdeq \trunc{2n}{\tsm{x:X} (\merid(x) \ct \opp{\merid(x_0)} = p)}\label{eq:freudcodeN}\\
    \code(\south,q) &\defeq \trunc{2n}{\hfib{\merid}{q}}
    \jdeq \trunc{2n}{\tsm{x:X} (\merid(x) = q)}.\label{eq:freudcodeS}
  \end{align}
\end{defn}

我们的最终目标是证明对所有 $y:\susp X$ 和 $p:\north=y$，$\code(y,p)$ 都是可缩的。
取 $y\defeq \north$，这将表明 $\sigma$ 的所有纤维都是 $2n$-连通的，从而 $\sigma$ 是 $2n$-连通的。

\begin{proof}[\cref{thm:freudcode} 的证明]
  我们通过对 $y:\susp X$ 进行归纳来定义 $\code(y,p)$，其中前两种情形是~\eqref{eq:freudcodeN} 与~\eqref{eq:freudcodeS}。
  接下来需要对每个 $x_1:X$ 构造一个依值道路
  \[ \dpath{\lam{y}(\north=y)\to\type}{\merid(x_1)}{\code(\north)}{\code(\south)}. \]
  根据 \cref{thm:dpath-arrow}，这等价于给出一族道路
  \[ \prd{q:\north=\south} \code(\north)(\transfib{\lam{y}(\north=y)}{\opp{\merid(x_1)}}{q}) = \code(\south)(q). \]
  而由泛等性（univalence）及道路类型中的传输，这又等价于给出一族等价
  \[ \prd{q:\north=\south} \code(\north,q \ct \opp{\merid(x_1)}) \eqvsym \code(\south,q). \]
  我们将定义一族映射
  \begin{equation}\label{eq:freudmap}
    \prd{q:\north=\south} \code(\north,q \ct \opp{\merid(x_1)}) \to \code(\south,q).
  \end{equation}
  然后证明它们都是等价的。
  因此，令 $q:\north=\south$；根据截断的泛性质以及 $\code(\north,\blank)$ 和 $\code(\south,\blank)$ 的定义，只需对每个 $x_2:X$ 定义一个映射
  \begin{equation*}
    \big(\merid(x_2)\ct \opp{\merid(x_0)} = q \ct \opp{\merid(x_1)}\big)
    \to \trunc{2n}{\tsm{x:X} (\merid(x) = q)}.
  \end{equation*}
  现在，对于每个 $x_1,x_2:X$，该类型是 $2n$-截断的，而 $X$ 是 $n$-连通的。
  于是，由 \cref{thm:wedge-connectivity}，只需分别在 $x_1$ 为 $x_0$ 和 $x_2$ 为 $x_0$ 的情形下定义此映射，并验证当两者同时为 $x_0$ 时它们一致即可。

  当 $x_1$ 为 $x_0$ 时，假设条件是 $r:\merid(x_2)\ct \opp{\merid(x_0)} = q \ct \opp{\merid(x_0)}$。
  由此，从 $r$ 中消去 $\opp{\merid(x_0)}$ 得到 $r':\merid(x_2)=q$，于是我们可以将其像定义为 $\tproj{2n}{(x_2,r')}$。

  当 $x_2$ 为 $x_0$ 时，假设条件是 $r:\merid(x_0)\ct \opp{\merid(x_0)} = q \ct \opp{\merid(x_1)}$。
  将其重新排列，我们得到 $r'':\merid(x_1)=q$，从而可以将其像定义为 $\tproj{2n}{(x_1,r'')}$。

  最后，当 $x_1$ 和 $x_2$ 均为 $x_0$ 时，只需证明得到的 $r'$ 和 $r''$ 是一致的；这是一个关于道路复合的简单引理。
  至此，我们完成了~\eqref{eq:freudmap} 的定义。
  为了证明它是一族等价，由于作为等价这一性质是一个命题（mere proposition）且 $x_0:\unit\to X$ 是（至少）$(-1)$-连通的，只需假设 $x_1$ 是 $x_0$。
  在这种情况下，检查上述构造可知，它本质上是消去 $\opp{\merid(x_0)}$ 这一函数的 $2n$-截断，而该函数本身是一个等价。
\end{proof}

除了~\eqref{eq:freudcodeN} 与~\eqref{eq:freudcodeS} 之外，我们还需要从 $\code$ 的构造中提取一些关于它如何作用于道路的信息。
为此，我们使用以下引理。

\begin{lem}\label{thm:freudlemma}
  设 $A:\UU$，$B:A\to \UU$，以及 $C:\prd{a:A} B(a)\to\UU$，同时给定 $a_1,a_2:A$ 以及 $m:a_1=a_2$ 和 $b:B(a_2)$。
  那么，函数
  \[\transfib{\widehat{C}}{\pairpath(m,t)}{\blank} : C(a_1,\transfib{B}{\opp m}{b}) \to C(a_2,b),\]
  （其中 $t:\transfib{B}{m}{\transfib{B}{\opp m}{b}} = b$ 是明显的相干道路，而 $\widehat{C}:(\sm{a:A} B(a)) \to\type$ 是 $C$ 的非柯里化形式）等于通过泛等性从以下复合得到的等价：
  \begin{align}
    C(a_1,\transfib{B}{\opp m}{b})
    &= \transfib{\lam{a} B(a)\to \UU}{m}{C(a_1)}(b)
    \tag{by~\eqref{eq:transport-arrow}}\\
    &= C(a_2,b). \tag{by $\happly(\apd{C}{m},b)$}
  \end{align}
\end{lem}

\begin{proof}
  根据道路归纳，我们可以假设 $a_2$ 就是 $a_1$ 且 $m$ 就是 $\refl{a_1}$，此时两个函数都是恒等函数。
\end{proof}

我们将此引理应用于以下情形：令 $A\defeq\susp X$，$B\defeq \lam{y}(\north=y)$，$C\defeq\code$；同时令 $a_1\defeq\north$，$a_2\defeq\south$，并对某个 $x_1:X$ 令 $m\defeq \merid(x_1)$；最后令 $b\defeq q$，即某条道路 $\north=\south$。
对 $\susp X$ 进行归纳的计算规则将 $\apd{C}{m}$ 等同于一条由泛等性和函数外延性以特定方式构造出的道路。
\cref{thm:freudlemma} 中描述的第二个函数，实质上就是撤销这些对泛等性和函数外延性的应用，从而归约到我们利用 \cref{thm:wedge-connectivity} 所定义的特定函数~\eqref{eq:freudmap}。
因此，\cref{thm:freudlemma} 表明，沿着 $\pairpath(q,t)$ 进行传输，实质上就恢复了这些函数。

最后，根据构造，当 $x_1$ 或 $x_2$ 与 $x_0$ 重合，并且输入在 $\tproj{2n}{\blank}$ 的像中时，我们能更明确地知道这些函数是什么。
于是，对任意 $x_2:X$，我们有
\begin{equation}
  \transfib{\hat{\code}}{\pairpath(\merid(x_0),t)}{\tproj{2n}{(x_2,r)}}
  =\tproj{2n}{(x_1,r')}\label{eq:freudcompute1}
\end{equation}
其中 $r:\merid(x_2) \ct \opp{\merid(x_0)} = \transfib{B}{\opp{\merid(x_0)}}{q}$ 如前所述是任意的，而 $r':\merid(x_2)=q$ 是由 $r$ 将其终点与 $q \ct \opp{\merid(x_0)}$ 等同并消去 $\opp{\merid(x_0)}$ 而得到的。
类似地，对任意 $x_1:X$，我们有
\begin{equation}
  \transfib{\hat{\code}}{\pairpath(\merid(x_1),t)}{\tproj{2n}{(x_0,r)}}
  = \tproj{2n}{(x_1,r'')}\label{eq:freudcompute2}
\end{equation}
其中 $r:\merid(x_0) \ct \opp{\merid(x_0)} = \transfib{B}{\opp{\merid(x_1)}}{q}$，而 $r'':\merid(x_1)=q$ 是通过将其终点等同起来并重新排列道路而得到的。

\begin{proof}[\cref{thm:freudenthal} 的证明]
  接下来要证明，对于每个 $y:\susp X$ 和 $p:\north=y$， $\code(y,p)$ 都是可缩的。
  首先，我们必须选定一个收缩中心，记作 $c(y,p):\code(y,p)$。
  这对应于我们先前证明中 $\encode$ 函数的定义，因此我们用传输来定义它。
  注意，在 $y$ 取 $\north$ 且 $p$ 取 $\refl{\north}$ 的特殊情况下，我们有
  \[\code(\north,\refl{\north}) \jdeq \trunc{2n}{\tsm{x:X} (\merid(x) \ct \opp{\merid(x_0)} = \refl{\north})}.\]
  因此，我们可以定义 $c(\north,\refl{\north})\defeq \tproj{2n}{(x_0,\mathsf{rinv}_{\merid(x_0)})}$，其中对于任意道路 $q$，$\mathrm{rinv}_q$ 是明显的道路 $q\ct\opp q = \refl{}$。
  接下来，我们可以通过对 $p$ 做道路归纳来得到 $c:\prd{y:\susp X}{p:\north=y} \code(y,p)$，但下面我们会看到，用传输给出一个具体的定义也很重要：
  \[ c(y,p) \defeq \transfib{\hat{\code}}{\pairpath(p,\mathsf{tid}_p)}{c(\north,\refl{\north})}
  \]
  这里 $\hat{\code}: \big(\sm{y:\susp X} (\north=y)\big) \to \type$ 是 \code 的非柯里化形式，而 $\mathsf{tid}_p:\trans{p}{\refl{}} = p$ 是一个标准引理。

  接着，我们必须证明 $\code(y,p)$ 中的每个元素都等于 $c(y,p)$。
  再次利用道路归纳，我们只需考虑 $y$ 为 $\north$ 且 $p$ 为 $\refl{\north}$ 的情形。
  事实上，我们将在更一般的情形下证明它，即仅假设 $y$ 是 $\north$，而 $p$ 是任意的。
  换言之，我们将证明，对于任意 $p:\north=\north$ 和 $d:\code(\north,p)$，都有 $d = c(\north,p)$。
  由于该等式是一个 $(2n-1)$-类型，我们可以假设 $d$ 具有 $\tproj{2n}{(x_1,r)}$ 的形式，其中 $x_1:X$，且 $r:\merid(x_1) \ct \opp{\merid(x_0)} = p$。

  现在，通过进一步的道路归纳，我们可以假定 $r$ 是自反性，并且 $p$ 是 $\merid(x_1) \ct \opp{\merid(x_0)}$。
  （这就是为什么我们在前面要推广到任意的 $p$。）
  于是，我们需要证明
  \begin{equation}
    \tproj{2n}{(x_1, \refl{\merid(x_1) \ct \opp{\merid(x_0)}})}
    \;=\;
    c\left(\north,\refl{\merid(x_1) \ct \opp{\merid(x_0)}}\right).\label{eq:freudgoal}
  \end{equation}
  根据定义，该等式右边是
  \begin{multline*}
    \Transfib{\hat{\code}}{\pairpath(\merid(x_1) \ct \opp{\merid(x_0)}, \nameless)}{\tproj{2n}{(x_0,\nameless)}} \\
    = \transfibf{\hat{\code}}
    \begin{aligned}[t]
      \Big(
      &{\pairpath(\opp{\merid(x_0)}, \nameless)},\\
      &{\Transfib{\hat{\code}}{\pairpath(\merid(x_1), \nameless)}{\tproj{2n}{(x_0,\nameless)}}}
      \Big)
    \end{aligned}
    \\
    = \Transfib{\hat{\code}}{\pairpath(\opp{\merid(x_0)}, \nameless)}{\tproj{2n}{(x_1,\nameless)}}
    = \tproj{2n}{(x_1,\nameless)}
  \end{multline*}
  其中下划线 $\nameless$ 处应填入合适的相干道路。
  这里，第一步使用了传输的函子性，第二步调用了~\eqref{eq:freudcompute2}，第三步则调用了~\eqref{eq:freudcompute1}（同时将传输移到等式的另一边）。
  这样一来，等式右边的第一分量就与~\eqref{eq:freudgoal} 左边的第一分量相同了。
  我们留给读者验证：所有相干道路都会相消，从而在第二分量中给出自反性。
\end{proof}

% As a corollary, we have the following equivalence.

\begin{cor}[Freudenthal 等价性] \label{cor:freudenthal-equiv}
  设 $X$ 是 $n$-连通且带点的类型，其中 $n\geq 0$。
  那么 $\eqv{\trunc{2n}{X}}{\trunc{2n}{\Omega\susp(X)}}$。
\end{cor}

\begin{proof}
  根据 \cref{thm:freudenthal}，$\sigma$ 是 $2n$-连通的。再由
  \cref{lem:connected-map-equiv-truncation} 可知，它因此在 $2n$-截断上诱导一个等价。
\end{proof}

\index{encode-decode method|)}%

\index{Freudenthal suspension theorem|)}%
\index{theorem!Freudenthal suspension|)}%

\index{homotopy!group!of sphere}%
\index{stability!of homotopy groups of spheres}%
\index{type!n-sphere@$n$-sphere}%
Freudenthal 纬悬定理的一个重要推论是，球面的同伦群在某个范围内是稳定的（这些范围对应于 \cref{tab:homotopy-groups-of-spheres} 中从东北到西南方向的对角线）：

\begin{cor}[球面的稳定性] \label{cor:stability-spheres}
  若 $k \le 2n-2$，则 $\pi_{k+1}(S^{n+1}) = \pi_{k}(S^{n})$。
\end{cor}

\begin{proof}
  假设 $k \le 2n-2$。
  %
  根据 \cref{cor:sn-connected}，$\Sn ^{n}$ 是 $\nminusone$-连通的。因此，
  由 \cref{cor:freudenthal-equiv} 可得，
  \[
\trunc{2(n-1)}{\Omega(\susp(\Sn^{n}))} = \trunc{2(n-1)}{\Sn^{n}}.
\]
  再根据 \cref{lem:truncation-le}，由于 $k \le 2(n-1)$，在等式两边同时应用 $\trunc{k}{\blank}$ 可知此式对 $k$ 也成立：
  \begin{equation}\label{eq:freudenthal-for-spheres}
\trunc{k}{\Omega(\susp(\Sn^{n}))} = \trunc{k}{\Sn^{n}}.
\end{equation}
  %
  于是，证明的主要想法如下；我们略去以下验证：这些等价在相应空间的基点上是否作用恰当，以及当 $k > 0$ 时这些等价是否保持乘法结构：
  %
  \begin{align*}
\pi_{k+1}(\Sn^{n+1}) &\jdeq \trunc{0}{\Omega^{k+1}(\Sn^{n+1})} \\
                     &\jdeq \trunc{0}{\Omega^k(\Omega(\Sn^{n+1}))} \\
                     &\jdeq \trunc{0}{\Omega^k(\Omega(\susp(\Sn^{n})))} \\
                     &= \Omega^k(\trunc{k}{(\Omega(\susp(\Sn^{n})))})
                     \tag{by \cref{thm:path-truncation}}\\
                     &= \Omega^k(\trunc{k}{\Sn^{n}})
                     \tag{by \eqref{eq:freudenthal-for-spheres}}\\
                     &= \trunc{0}{\Omega^k(\Sn^{n})}
                     \tag{by \cref{thm:path-truncation}}\\
                     &\jdeq \pi_k(\Sn^{n}). \qedhere
\end{align*}
  %
\end{proof}

这意味着，一旦我们计算出了某条稳定对角线上的一个条目，我们就知道了该对角线上所有的条目。例如：

\begin{thm}\label{thm:pinsn}
对每个 $n\geq 1$，有 $\pi_n(\Sn^n)=\Z$。
\end{thm}

\begin{proof}
证明对 $n$ 施行归纳。我们已知 $\pi_1(\Sn ^1) = \Z$（\cref{cor:pi1s1}）与 $\pi_2(\Sn ^2) = \Z$（\cref{cor:pis2-hopf}）。
当 $n \ge 2$ 时，$n \le (2n - 2)$。于是，由
\cref{cor:stability-spheres}、$\pi_{n+1}(S^{n+1}) = \pi_{n}(S^{n})$ 以及这一等价，再结合归纳假设，即得结论。
\end{proof}

\begin{cor}
对任意 $n\ge -1$，$\Sn^{n+1}$ 不是 $n$-类型。
\end{cor}

\begin{cor}\label{thm:pi3s2}
$\pi_3(\Sn^2)=\Z$。
\end{cor}

\begin{proof}
由 \cref{cor:pis2-hopf}，$\pi_3(\Sn^2) = \pi_3(\Sn^3)$。
而根据 \cref{thm:pinsn}，有 $\pi_3(\Sn^3) = \Z$。
\end{proof}

\section{van Kampen 定理}
\label{sec:van-kampen}

\index{van Kampen theorem|(}%
\index{theorem!van Kampen|(}%

\index{fundamental!group}%
van Kampen 定理用于计算空间的（同伦）推出的基本群 $\pi_1$。
传统上，这一定理针对一个拓扑空间 $X$，它是两个开子空间 $U$ 和 $V$ 的并集；
但在同伦论的意义下，这只是确保 $X$ 是 $U$ 和 $V$ 沿其交集的推出的一种简便方式。
因此，我们将证明一个适用于任意推出的 van Kampen 定理版本。

在本节中，我们将使用与 \cref{sec:pi1-s1-intro} 中计算 $\pi_1(\Sn^1)$ 时相同的编码-解码方法来证明 van Kampen 定理。
此外，也存在一种更偏向同伦论的处理方法；参见 \cref{ex:rezk-vankampen}。

我们需要编码-解码方法的一个更精细的版本。
在 \cref{sec:pi1-s1-intro}（以及 \cref{sec:compute-coprod,sec:compute-nat}）中，我们用它来刻画一个（高阶）归纳类型 $W$ 的道路空间——从而得到环路空间 $\Omega(W)$ 的刻画，并进而得到其 0-截断 $\pi_1(W)$ 的刻画。
在 van Kampen 定理中，我们的目标只是刻画基本群 $\pi_1(W)$，而我们没有任何环路空间或道路空间的显式描述可供使用。

事实证明，我们可以将同样的技术直接用于道路纤维化的截断版本，从而不仅刻画基本\emph{群} $\pi_1(W)$，还刻画整个基本\emph{群胚}。
\index{fundamental!pregroupoid}%
具体来说，对于一个类型 $X$，记 $\Pi_1 X: X\to X\to \type$ 为其相等类型的 $0$-截断，即$\Pi_1 X(x,y) \defeq \trunc0{x=y}$。
注意我们有诱导的群胚运算
\begin{align*}
  (\blank\ct\blank) &\;:\; \Pi_1X(x,y) \to \Pi_1X(y,z) \to \Pi_1X(x,z)\\
  \opp{(\blank)} &\;:\; \Pi_1X(x,y) \to \Pi_1X(y,x)\\
  \refl{x} &\;:\; \Pi_1X(x,x)\\
  \apfunc{f} &\;:\; \Pi_1X(x,y) \to \Pi_1Y(fx,fy)
\end{align*}
我们对其使用与道路上的相应运算相同的记号。

\subsection{朴素 van Kampen}
\label{sec:naive-vankampen}

我们从 van Kampen 定理的一个``朴素''版本开始，它虽有用，但不及经典版本。
在 \cref{sec:better-vankampen} 中，我们将把它改进为一个更有用的版本。

\index{encode-decode method|(}%

给定类型 $A,B,C$ 和函数 $f:A\to B$ 与 $g:A\to C$，令 $P$ 为它们的推出 $B\sqcup^A C$。
如我们在 \cref{sec:colimits} 中所见，$P$ 是由以下生成元生成的高阶归纳类型：


\begin{itemize}
\item $i:B\to P$，
\item $j:C\to P$，以及
\item 对所有 $x:A$，一条道路 $k x:ifx = jgx$。
\end{itemize}


通过对 $P$ 的双重归纳，定义 $\code:P\to P\to \type$ 如下。


\begin{itemize}
\item $\code(ib,ib')$ 是以下序列类型的一个集合商（见 \cref{sec:set-quotients}）：
  %pairs $(\vec a, \vec a', \vec p, \vec q)$
  \[ (b, p_0, x_1, q_1, y_1, p_1, x_2, q_2, y_2, p_2, \dots, y_n, p_n, b') \]
  其中
  \begin{itemize}
  \item $n:\mathbb{N}$
  \item 对 $0<k \le n$，有 $x_k:A$ 和 $y_k:A$
  \item $p_0:\Pi_1B(b,f x_1)$ 与 $p_n:\Pi_1B(f y_n, b')$（当 $n>0$），以及 $p_0:\Pi_1B(b,b')$（当 $n=0$）
  \item $p_k:\Pi_1B(f y_k, fx_{k+1})$（当 $1\le k < n$）
  \item $q_k:\Pi_1C(gx_k, gy_k)$（当 $1\le k\le n$）
  \end{itemize}
  该商由以下等式生成：
  \begin{align*}
    (\dots, q_k, y_k, \refl{fy_k}, y_k, q_{k+1},\dots)
    &= (\dots, q_k\ct q_{k+1},\dots)\\
    (\dots, p_k, x_k, \refl{gx_k}, x_k, p_{k+1},\dots)
    &= (\dots, p_k\ct p_{k+1},\dots)
  \end{align*}
  （见下文 \cref{rmk:naive}）。
  我们留给读者将这种序列类型精确定义为一个归纳类型。
\item $\code(jc,jc')$ 的定义类似，但交换 $B$ 和 $C$ 的角色。
  我们同样在记号上交换 $x$ 和 $y$ 以及 $p$ 和 $q$ 的角色。
\item $\code(ib,jc)$ 和 $\code(jc,ib)$ 的定义也类似，但改变奇偶性，使其起始于一个类型而终止于另一个类型。
\item 对于 $a:A$ 和 $b:B$，我们需要一个等价
  \begin{equation}
    \code(ib, ifa) \eqvsym \code(ib,jga).\label{eq:bfa-bga}
  \end{equation}
  我们将其定义为由如下定义在序列上的两个函数给出：
  \begin{align*}
    (\dots, y_n, p_n,fa) &\mapsto (\dots,y_n,p_n,a,\refl{ga},ga),\\
    (\dots, x_n, p_n, a, \refl{fa}, fa) &\mapsfrom (\dots, x_n, p_n, ga).
  \end{align*}
  这两个函数都容易看出保持等价关系，从而在代码类型上定义了函数。
  从左到右再回到左的复合为
  \[ (\dots, y_n, p_n,fa) \mapsto
  (\dots,y_n,p_n,a,\refl{ga},a,\refl{fa},fa)
  \]
  根据商的一个生成等式，它等于恒等映射。
  另一个方向的复合类似。
  这样我们就定义了一个等价~\eqref{eq:bfa-bga}。
\item 类似地，我们需要等价
  \begin{align*}
    \code(jc,ifa) &\eqvsym \code(jc,jga)\\
    \code(ifa,ib)&\eqvsym (jga,ib)\\
    \code(ifa,jc)&\eqvsym (jga,jc)
  \end{align*}
  它们都以完全相同的方式定义（后两个是通过在序列开头而非结尾添加自反性项来定义的）。
\item 最后，我们需要知道，对于 $a,a':A$，下图交换：
  \begin{equation}\label{eq:bfa-bga-comm}
  \vcenter{\xymatrix{
      \code(ifa,ifa') \ar[r]\ar[d] &
      \code(ifa,jga')\ar[d]\\
      \code(jga,ifa')\ar[r] &
      \code(jga,jga')
      }}
  \end{equation}
  这相当于说，如果我们先在一个序列的开头添加一些东西，然后在结尾再添加一些东西，那么换一种顺序来做也是一样的。
\end{itemize}

\begin{rmk}\label{rmk:naive}
  人们可能期望在 \code 的定义中看到集合商的一些额外的生成等式，例如
  \begin{align*}
    (\dots, p_{k-1} \ct fw, x_{k}', q_{k}, \dots) &=
    (\dots, p_{k-1}, x_{k}, gw \ct q_{k}, \dots)
    \tag{for $w:\Pi_1A(x_{k},x_{k}')$}\\
    (\dots, q_k \ct gw, y_k', p_k, \dots) &=
    (\dots, q_k, y_k, fw \ct p_k, \dots).
    \tag{for $w:\Pi_1A(y_k, y_k')$}
  \end{align*}
  然而，这些并非必需！
  事实上，它们可以通过对 $w$ 进行道路归纳自动得出。
  这正是``朴素''van Kampen 定理与下一小节将要考虑的更精细版本之间的主要区别。
\end{rmk}

继续下去，我们可以刻画在纤维化 $\code$ 中的 transport 操作：


\begin{itemize}
\item 对于 $p:b=_B b'$ 和 $u:P$，我们有
  \[ \mathsf{transport}^{b\mapsto \code(u,ib)}(p, (\dots, y_n,p_n,b))
  = (\dots,y_n,p_n\ct p,b').
  \]
\item 对于 $q:c=_C c'$ 和 $u:P$，我们有
  \[ \mathsf{transport}^{c\mapsto \code(u,jc)}(q, (\dots, x_n,q_n,c))
  = (\dots,x_n,q_n\ct q,c').
  \]
\end{itemize}


这里我们在记号上有所滥用：对 $X$ 中的一条道路及其在 $\Pi_1X$ 中的像使用了相同的名称。
注意，$\Pi_1X$ 中的 transport 也是通过与一条道路（的像）进行拼接给出的。
由此，我们可以通过对 $u$ 进行归纳来证明上述命题。
我们还有：


\begin{itemize}
\item 对于 $a:A$ 和 $u:P$，
  \[ \mathsf{transport}^{v\mapsto \code(u,v)}(ha, (\dots, y_n,p_n,fa))
  = (\dots,y_n,p_n,a,\refl{ga},ga).
  \]
\end{itemize}


这本质上由 $\code$ 的定义得出。

我们还要通过对 $u$ 进行归纳来构造一个函数
\[ r : \prd{u:P} \code(u,u) \]
具体如下：
\begin{align*}
  rib &\defeq (b,\refl{b},b)\\
  rjc &\defeq (c,\refl{c},c)
\end{align*}
而对于 $rka$，我们取复合相等
\begin{align*}
  (ka,ka)_* (fa,\refl{fa},fa)
  &= (ga,\refl{ga},a,\refl{fa},a,\refl{ga},ga) \\
  &= (ga,\refl{ga},ga)
\end{align*}
其中第一个相等来自上述关于在 $\code$ 中 transport 的观察，第二个相等则是定义 $\code$ 时所使用的集合商关系的一个实例。

接下来我们将证明：


\begin{thm}[朴素 van Kampen 定理]\label{thm:naive-van-kampen}
  对于所有 $u,v:P$，存在一个等价
  \[ \Pi_1P(u,v) \eqvsym \code(u,v). \]
\end{thm}

\begin{proof}

为定义一个函数
\[ \encode : \Pi_1P(u,v) \to \code(u,v) \]
只需定义一个函数 $(u=_P v) \to \code(u,v)$ 即可，因为 $\code(u,v)$ 是一个集合。
我们通过 transport 来定义它：
\[\encode(p) \defeq \mathsf{transport}^{v\mapsto \code(u,v)}(p,r(u)).\]
现在来定义
\[ \decode: \code(u,v) \to \Pi_1P(u,v) \]
我们照常对 $u,v:P$ 进行归纳。
在 $u$ 和 $v$ 的每一种情形中，我们将 $i$ 或 $j$ 适当地作用于所有等式 $p_k$ 和 $q_k$，并在 $P$ 中拼接结果，同时利用 $h$ 来等同端点。
例如，当 $u\jdeq ib$ 且 $v\jdeq ib'$ 时，我们定义
\begin{narrowmultline}\label{eq:decode}
 \decode(b, p_0, x_1, q_1, y_1, p_1, \dots, y_n, p_n, b') \defeq\narrowbreak
 (p_0)\ct h(x_1) \ct j(q_1) \ct \opp{h(y_1)} \ct i(p_1) \ct \cdots \ct \opp{h(y_n)}\ct i(p_n).
\end{narrowmultline}
这个定义保持集合商的等价关系以及诸如 \eqref{eq:bfa-bga} 的等价，因为 $h: fi \htpy gj$ 是自然的，且 $f$ 和 $g$ 具有函子性。

照例，为了证明复合
\[ \Pi_1P(u,v) \xrightarrow{\encode} \code(u,v) \xrightarrow{\decode} \Pi_1P(u,v) \]
是恒等映射，我们首先``剥掉''0-截断（因为值域是集合），然后应用道路归纳。
输入 $\refl{u}$ 映到 $ru$，而 $ru$ 又映回 $\refl u$（通过对 $u$ 的进一步归纳来分解 $\decode(ru)$）。

最后，考虑复合
\[  \code(u,v) \xrightarrow{\decode} \Pi_1P(u,v) \xrightarrow{\encode} \code(u,v). \]
我们对 $u,v:P$ 进行归纳。
当 $u\jdeq ib$ 且 $v\jdeq ib'$ 时，此复合是
%
\begin{narrowmultline*}
(b, p_0, x_1, q_1, y_1, p_1, \dots, y_n, p_n, b')
\narrowbreak
\begin{aligned}[t]
  &\mapsto \Big(ip_0\ct hx_1 \ct jq_1 \ct \opp{hy_1} \ct ip_1 \ct \cdots \ct \opp{hy_n}\ct ip_n\Big)_*(rib)\\
  &= (ip_n)_* \cdots(jq_1)_* (hx_1)_*(ip_0)_*(b,\refl{b},b)\\
  &= (ip_n)_* \cdots(jq_1)_* (hx_1)_*(b,p_0,ifx_1)\\
  &= (ip_n)_* \cdots(jq_1)_* (b,p_0,x_1,\refl{gx_1},jgx_1)\\
  &= (ip_n)_* \cdots (b,p_0,x_1,q_1,jgy_1)\\
  &= \quad\vdots\\
  &= (b, p_0, x_1, q_1, y_1, p_1, \dots, y_n, p_n, b').
\end{aligned}
\end{narrowmultline*}
%
即恒等函数。
（准确地说，这里需要一个隐式的归纳论证。）
另外三种点的情形是类似的，而道路的情形则是平凡的，因为所涉及的类型都是集合。
\end{proof}

\index{encode-decode method|)}%

\cref{thm:naive-van-kampen} 使我们能够计算许多类型的基本群，前提是 $A$ 是一个集合；因为在这种情况下，根据定义，每个 $\code(u,v)$ 都是一个\emph{集合}模掉一个关系所得的集合商。

\begin{eg}\label{eg:circle}
  令 $A\defeq \bool$，$B\defeq\unit$，且 $C\defeq \unit$。
  那么 $P \eqvsym S^1$。
  考察比如 $\code(i(\ttt),i(\ttt))$ 的定义，我们发现所有的道路实际上都是平凡的，因此唯一的信息就是元素的序列 $x_1,y_1,\dots,x_n,y_n: \bool$。
  此外，若对某个 $k$ 我们有 $x_k=y_k$ 或 $y_k=x_{k+1}$，那么集合商关系允许我们同时``切除''这两个元素。
  因此，每个这样的序列都等于一个典范的\emph{约化}序列，其中没有两个相邻元素是相等的。
  显然，这样的约化序列由其长度（一个自然数 $n$）以及（如果 $n>1$）$x_1$ 是 $\bfalse$ 还是 $\btrue$ 这一信息所唯一决定，因为这能唯一确定序列的其余部分。
  当然，这些数据可以等同于一个整数：$n$ 对应绝对值，$x_1$ 编码正负号。
  于是我们重新得到了 $\pi_1(S^1)\cong \Z$。
\end{eg}

由于 \cref{thm:naive-van-kampen} 断言的只是集合族之间的双射，这个同构 $\pi_1(S^1)\cong \Z$ 同样也只是集合间的双射。
然而，我们可以在 $\code$ 上定义一种拼接运算（通过拼接序列），并证明 $\encode$ 和 $\decode$ 构成了一个保持该结构的同构。
（用 \cref{cha:category-theory} 的语言来说，这些将是``预群胚''。）
细节留给读者。

\index{fundamental!group|(}%

\begin{eg}\label{eg:suspension}
  更一般地，令 $B\defeq\unit$ 且 $C\defeq \unit$，但 $A$ 是任意的，于是 $P$ 是 $A$ 的纬悬。
  此时道路 $p_k$ 和 $q_k$ 也是平凡的，因此一个道路编码中唯一的信息是元素的序列 $x_1,y_1,\dots,x_n,y_n: A$。
  前两个生成等式表明相邻的相等元素可以消去，因此将这个序列视为如下形式的群中的字是合理的：
  \[ x_1 y_1^{-1} x_2 y_2^{-1} \cdots x_n y_n^{-1} \]
  实际上，它看起来类似于 $A$（或者说 $\trunc0A$；参见 \cref{thm:freegroup-nonset}）上的自由群，但我们只考虑那些以非逆元素开头、交替出现逆与非逆元素、并以逆元素结尾的字。
  这实际上将生成元集合的规模减少了一个。
  例如，如果 $A$ 有一个点 $a:A$，那么我们可以将 $\pi_1(\susp A)$ 等同于由 $\trunc0A$ 作为生成元且满足关系 $\tproj0a = e$ 所给出的群；详见 \cref{ex:vksusppt,ex:vksuspnopt}。
\end{eg}

\begin{eg}\label{eg:wedge}
  令 $A\defeq\unit$，且 $B$ 和 $C$ 是任意的，于是 $f$ 和 $g$ 只是给 $B$ 和 $C$ 配备了基点，比如 $b$ 和 $c$。
  那么 $P$ 就是 $B$ 和 $C$ 的\emph{楔} $B\vee C$（带基点空间范畴中的余积）。
  在这种情况下，元素 $x_k$ 和 $y_k$ 是平凡的，因此唯一的信息是一系列环路 $(p_0,q_1,p_1,\dots,p_n)$，其中 $p_k:\pi_1(B,b)$ 且 $q_k:\pi_1(C,c)$。
  这样的序列模去我们所施加的等价关系后，很容易与群 $\pi_1(B,b)$ 和 $\pi_1(C,c)$ 的\emph{自由积}的显式描述等同，如 \cref{sec:free-algebras} 中所构造的那样。
  因此，我们有 $\pi_1(B\vee C) \cong \pi_1(B) * \pi_1(C)$。
\end{eg}

\index{fundamental!group|)}%

然而，\cref{thm:naive-van-kampen} 离完整的经典 van Kampen 定理还差一步。经典版本能处理 $A$ 不一定是集合的情形，并且断言 $\pi_1(B\sqcup^A C) \cong \pi_1(B) *_{\pi_1(A)} \pi_1(C)$（其中的基点取自 $A$）。
事实上，\cref{thm:naive-van-kampen} 的结论完全没有涉及 $\pi_1(A)$；$A$ 中的道路以一种类型论的方式被``嵌入''到商化构造当中，这使得提取显式信息变得困难，因为 $\code(u,v)$ 是对一个非集合按某个关系做集合商得到的结果。
因此，在下一小节中，我们将考虑一个更好的 van Kampen 定理版本。

\subsection{带基点集的 van Kampen 定理}
\label{sec:better-vankampen}

\index{basepoint!set of}%
我们将要介绍的 van Kampen 定理的改进，与经典代数拓扑中的类似改进密切相关，即在经典理论中为 $A$ 配备一个\emph{基点集 $S$}。
事实上，我们的证明并不需要假设这个``基点集''是一个\emph{集合}——它完全可以是一个任意类型；假设 $S$ 是集合的用处要到之后应用定理进行计算时才体现出来。
重要的是，$S$ 包含 $A$ 的每个连通分支中至少一个点。
我们在类型论中将其表述为：存在一个类型 $S$ 和一个函数 $k:S \to A$，它是满的，亦即 $(-1)$-连通的。
如果 $S\jdeq A$ 且 $k$ 是恒等函数，那么我们将恢复朴素的 van Kampen 定理。
另一个值得记住的例子是，当 $A$ 是带点且 (0-)连通的，并令 $k:\unit\to A$ 为取基点的映射时：由 \cref{thm:minusoneconn-surjective,thm:connected-pointed}，此映射是满的恰好当 $A$ 是 0-连通的。

令 $A,B,C,f,g,P,i,j,h$ 如前一小节所述。
现在，给定满射 $k:S\to A$，我们定义一个辅助类型，用以改善 $k$ 的连通性。
令 $T$ 为由以下生成元构造的高阶归纳类型：

\begin{itemize}
\item 一个函数 $\ell:S\to T$，以及
\item 对每个 $s,s':S$，一个函数 $m:(\id[A]{ks}{ks'}) \to (\id[T]{\ell s}{\ell s'})$。
\end{itemize}

存在一个明显的诱导函数 $\kbar:T\to A$ 使得 $\kbar \ell = k$，并且任何 $p:ks=ks'$ 都等于复合 $ks = \kbar \ell s \overset{\kbar m p}{=} \kbar \ell s' = k s'$。

\begin{lem}\label{thm:kbar}
  $\kbar$ 是 0-连通的。
\end{lem}

\begin{proof}
  我们需要证明，对任意 $a:A$，类型 $\sm{t:T}(\kbar t = a)$ 的 0-截断是可缩的。
  由于可缩性是一个命题，而 $k$ 是 $(-1)$-连通的，我们不妨设 $a=ks$ 对某个 $s:S$ 成立。
  于是我们可以将收缩中心取为 $\tproj0{(\ell s,q)}$，其中 $q$ 是相等 $\kbar\ell s = k s$。

  剩下还需要证明，对于任意 $\phi:\trunc0{\sm{t:T} (\kbar t = ks)}$，我们有 $\phi = \tproj0{(\ell s,q)}$。
  因为后者是一个命题，特别地也是集合，我们可以进一步假设 $\phi=\tproj0{(t,p)}$ 对某个 $t:T$ 和 $p:\kbar t = ks$ 成立。

  现在我们可以对 $t:T$ 进行归纳。
  如果 $t\jdeq\ell s'$，那么由 $ks' = \kbar \ell s' \overset{p}{=} ks$ 可以通过 $m$ 导出一个相等 $\ell s = \ell s'$。
  于是，根据 $\kbar$ 的定义以及同伦纤维中相等的定义，我们得到相等 $(ks',p) = (ks,q)$，从而得到 $\tproj0{(ks',p)} = \tproj0{(ks,q)}$。
  接下来我们必须证明，当 $t$ 沿着 $m$ 变化时，这些相等彼此相容。
  但这些相等是在一个集合（即 $\trunc0{\sm{t:T} (\kbar t = ks)}$）中的，因此这是自动的。
\end{proof}

\begin{rmk}
  \index{kernel!pair}%
  $T$ 可以被视为 $k$ 的``核偶对''的（同伦）余等化子。
  如果 $S$ 和 $A$ 是集合，那么 $k$ 的 $(-1)$-连通性将蕴含 $A$ 就是这个余等化子的 $0$-截断（参见 \cref{cha:set-math}）。
  对于一般的类型，高阶意象理论则提示，$k$ 的 $(-1)$-连通性将转而蕴含 $A$ 是 $k$ 的``单纯核''的余极限（亦称``几何实现''）。
  \index{.infinity1-topos@$(\infty,1)$-topos}%
  \index{geometric realization}%
  \index{simplicial!kernel}%
  \index{kernel!simplicial}%
  类型 $T$ 是这个单纯核的``1-骨架''的余极限，因此，它能将 $k$ 的连通性提高 $1$ 是合理的。
  更一般地，我们可以期望 $n$-骨架\index{skeleton!of a CW-complex}的余极限将连通性提高 $n$。
\end{rmk}

\index{encode-decode method|(}%

现在，我们通过双重归纳如下定义 $\code:P\to P\to \type$：

\begin{itemize}
\item $\code(ib,ib')$ 现在是如下类型序列的集合商：
  \[ (b, p_0, x_1, q_1, y_1, p_1, x_2, q_2, y_2, p_2, \dots, y_n, p_n, b') \]
  其中
  \begin{itemize}
  \item $n:\mathbb{N}$，
  \item 对 $0<k \le n$，有 $x_k:S$ 和 $y_k:S$，
  \item 当 $n>0$ 时，$p_0:\Pi_1B(b,f k x_1)$ 且 $p_n:\Pi_1B(f k y_n, b')$；当 $n=0$ 时，$p_0:\Pi_1B(b,b')$，
  \item 对 $1\le k < n$，有 $p_k:\Pi_1B(fk y_k, fkx_{k+1})$，
  \item 对 $1\le k\le n$，有 $q_k:\Pi_1C(gkx_k, gky_k)$。
  \end{itemize}
  该商由以下相等生成（参见\cref{rmk:naive}）：
  \begin{align*}
    (\dots, q_k, y_k, \refl{fy_k}, y_k, q_{k+1},\dots)
    &= (\dots, q_k\ct q_{k+1},\dots)\\
    (\dots, p_k, x_k, \refl{gx_k}, x_k, p_{k+1},\dots)
    &= (\dots, p_k\ct p_{k+1},\dots)\\
    (\dots, p_{k-1} \ct fw, x_{k}', q_{k}, \dots) &=
    (\dots, p_{k-1}, x_{k}, gw \ct q_{k}, \dots)
    \tag{for $w:\Pi_1A(kx_{k},kx_{k}')$}\\
    (\dots, q_k \ct gw, y_k', p_k, \dots) &=
    (\dots, q_k, y_k, fw \ct p_k, \dots).
    \tag{for $w:\Pi_1A(ky_k, ky_k')$}
  \end{align*}
  我们后面会用到在此类序列上 $\decode$ 的定义，它和之前一样，将所有路径 $p_k$ 和 $q_k$ 与 $h$ 的实例拼接起来，以得到 $\Pi_1P(ifb,ifb')$ 中的一个元素，参见~\eqref{eq:decode}。
  和之前一样，其他三种点的情况几乎相同。
\item 对于 $a:A$ 和 $b:B$，我们需要一个等价：
  \begin{equation}
    \code(ib, ifa) \eqvsym \code(ib,jga).\label{eq:bfa-bga2}
  \end{equation}
  由于 $\code$ 是集合值的，根据\cref{thm:kbar}，我们可以假设 $a=\kbar t$，其中 $t:T$。
  接下来，我们可以对 $t$ 进行归纳。
  若 $t\jdeq \ell s$，其中 $s:S$，那么我们如\cref{sec:naive-vankampen}那样定义~\eqref{eq:bfa-bga2}：
  \begin{align*}
    (\dots, y_n, p_n,fks) &\mapsto (\dots,y_n,p_n,s,\refl{gks},gks),\\
    (\dots, x_n, p_n, s, \refl{fks}, fks) &\mapsfrom (\dots, x_n, p_n, gks).
  \end{align*}
  这些定义尊重等价关系，并且如前所述定义了拟逆。
  现在假设 $t$ 沿着某个 $w:ks=ks'$ 对应的 $m_{s,s'}(w)$ 变化；我们必须证明~\eqref{eq:bfa-bga2} 尊重沿着 $\kbar mw$ 的传输。
  根据 $\kbar$ 的定义，这本质上可归结为沿 $w$ 本身的传输。
  由道路类型中传输的特征描述，我们需要证明的是：
  \[ w_*(\dots, y_n, p_n,fks) = (\dots,y_n, p_n \ct fw, fks') \]
  被~\eqref{eq:bfa-bga2} 映射到
  \[ w_*(\dots,y_n,p_n,s,\refl{gks},gks) = (\dots, y_n, p_n, s, \refl{gks} \ct gw, gks') \]
  但这直接由定义 \code 的集合商时所施加的新生成元得出。
\item 其他三种必需的等价以类似方式定义。
\item 最后，由于交换性~\eqref{eq:bfa-bga-comm} 是一个命题，根据 $k$ 的 $(-1)$-连通性，我们可以假设 $a=ks$ 且 $a'=ks'$，在这种情况下，其证明与之前完全相同。
\end{itemize}

\begin{thm}[带基点集的 van Kampen 定理]\label{thm:van-Kampen}
  对于所有 $u,v:P$，存在一个等价：
  \[ \Pi_1P(u,v) \eqvsym \code(u,v). \]
  其中 \code 的定义如本节所述。
\end{thm}

\begin{proof}
  基本上和之前一样。
  为了证明 $\decode$ 尊重商关系的新生成元，我们利用了 $h$ 的自然性。
  为了证明 $\decode$ 尊重诸如~\eqref{eq:bfa-bga2}之类的等价，我们需要对 $\kbar$ 和 $T$ 进行归纳，以便将这些等价分解为它们的定义式，但随后这又简单地变为 $f$ 和 $g$ 的函子性。
  其余部分是容易的。
  特别地，对于 $\encode\circ\decode$ 不需要额外的论证，因为目标是在集合中证明一个相等，所以关于 $h$ 的情况是平凡的。
\end{proof}

\index{encode-decode method|)}%

\index{fundamental!group|(}%
\cref{thm:van-Kampen} 允许我们计算空间~$A$ 的基本群，
即使 $A$ 不是集合，只要 $S$ 是集合；因为在此情况下，
根据定义，每个 $\code(u,v)$ 都是对一个\emph{集合}按某个关系做集合商的结果。
在这一点上，它是对
\cref{thm:naive-van-kampen}
的一个改进。

\begin{eg}\label{eg:clvk}
  假设 $S\defeq \unit$，那么 $A$ 有一个基点 $a \defeq k(\ttt)$ 并且是连通的。
  此时，推出的环路的编码可以等同于 $\pi_1(B,f(a))$ 和 $\pi_1(C,g(a))$ 中环路的交错序列，模去一个等价关系；该关系允许我们在它们之间滑动 $\pi_1(A,a)$ 的元素（分别在应用 $f$ 和 $g$ 之后）。
  因此，$\pi_1(P)$ 可以等同于 \emph{融合自由积}
  \index{amalgamated free product}%
  \index{free!product!amalgamated}%
  $\pi_1(B) *_{\pi_1(A)} \pi_1(C)$（群范畴中的推出），如\cref{sec:free-algebras}所构造。
  这（当 $B$ 和 $C$ 是 $P$ 的开子空间且 $A$ 是它们的交集时）大概是 van Kampen 定理最经典的版本。
\end{eg}

\begin{eg}\label{eg:cofiber}
  \index{cofiber}
  作为\cref{eg:clvk}的一个特例，我们再额外假设 $C\defeq\unit$，于是 $P$ 是余纤维 $B/A$。
  那么 $C$ 中的每个环路都等于自反性，所以道路编码上的关系允许我们将所有序列塌缩为 $B$ 中的单个环路。
  附加的关系要求在左边、右边或中间乘以 $\pi_1(A)$ 的像中的一个元素后等于单位元。
  因此，我们可以将 $\pi_1(B/A)$ 等同于群 $\pi_1(B)$ 模去由 $\pi_1(A)$ 的像生成的正规子群所得的商群。
\end{eg}

\begin{eg}\label{eg:torus}
  \index{torus}
  作为\cref{eg:cofiber}的进一步特例，令 $B\defeq S^1 \vee S^1$，令 $A\defeq S^1$，并令 $f:A\to B$ 选取复合环路 $p \ct q \ct \opp p \ct \opp q$，其中 $p$ 和 $q$ 是组成 $B$ 的两个 $S^1$ 拷贝中的生成环路。
  此时 $P$ 就是环面 $T^2$ 的一个表示。
  实际上，不难将 $P$ 与\cref{sec:hubs-spokes}中使用特定环路的锥所描述的 $T^2$ 表示等同起来。
  因此，$\pi_1(T^2)$ 是两个生成元上的自由群\index{generator!of a group}（即 $\pi_1(B)$）模去关系$p \ct q \ct \opp p \ct \opp q = 1$所得的商群。
  这显然给出两个生成元上的自由\emph{交换}群\index{group!abelian}，即 $\Z\times\Z$。
\end{eg}

\begin{eg}
  \index{CW complex}
  \index{hub and spoke}
  更一般地，任何 CW 复形都可以通过反复对球面“取锥”而得到，如\cref{sec:hubs-spokes}所述。
  具体来说，我们从一组点（“0-胞腔”）构成的集合 $X_0$ 出发，它称为 CW 复形的“0-骨架”。
  对某个 1-胞腔集合 $S_1$ 和某个“附着映射”族 $f_1$ 取推出
  \begin{equation*}
    \vcenter{\xymatrix{
        S_1 \times \Sn^0\ar[r]^-{f_1}\ar[d] &
        X_0\ar[d]\\
        \unit \ar[r] &
        X_1
      }}
  \end{equation*}
  得到“1-骨架”\index{skeleton!of a CW-complex} $X_1$。
  \index{attaching map}%
  然后再对某个 2-胞腔集合 $S_2$ 和某个附着映射族 $f_2$ 取推出
  \begin{equation*}
    \vcenter{\xymatrix{
        S_2 \times \Sn^1\ar[r]^{f_2}\ar[d] &
        X_1\ar[d]\\
        \unit \ar[r] &
        X_2
      }}
  \end{equation*}
  得到 2-骨架 $X_2$，依此类推。
  每次取推出所得的基本群都可以用 van Kampen 定理计算：得到一个群表示，其生成元由 1-骨架导出，关系由 $S_2$ 和 $f_2$ 导出。
  此后各步的推出不再改变基本群，因为当 $n>1$ 时 $\pi_1(\Sn^n)$ 是平凡的（参见\cref{sec:pik-le-n}）。
\end{eg}

\begin{eg}\label{eg:kg1}
  特别地，假设给定任意一个（集合意义上的）群 $G = \langle X \mid R \rangle$ 的表示\index{presentation!of a group}，其中 $X$ 是生成元集合，$R$ 是这些生成元\index{generator!of a group}上的字构成的集合。
  令 $B\defeq \bigvee_X S^1$ 且 $A\defeq \bigvee_R S^1$，并将 $f:A\to B$ 定义为将每个 $S^1$ 拷贝映到 $B$ 的生成环路中对应的字。
  由此可得 $\pi_1(P) \cong G$；这样我们就构造了一个基本群为 $G$ 的连通类型。
  由于任何群都有表示，所以任何群都可以成为某个类型的基本群。
  如果我们对这样一个类型取 1-截断，就得到一个类型，其唯一的非平凡同伦群为 $G$；这称为\define{艾伦伯格–麦克兰恩空间} $K(G,1)$。%
  \indexdef{Eilenberg--Mac Lane space}%
\end{eg}

\index{fundamental!group|)}%

\index{van Kampen theorem|)}%
\index{theorem!van Kampen|)}%

\section{Whitehead 定理与 Whitehead 原理}
\label{sec:whitehead}

在经典同伦论中，若映射 $f:A\to B$ 对 $A$ 中的所有点 $a$ 都诱导同构$\pi_n(A,a) \cong \pi_n(B,f(a))$（同时也诱导同构$\pi_0(A)\cong\pi_0(B)$），那么只要空间 $A$ 和 $B$ 性质良好（例如具有 CW 复形的同伦类型），$f$ 就必然是一个同伦等价。
\index{theorem!Whitehead's}%
\index{Whitehead's!theorem}%
这被称为\emph{Whitehead 定理}。
实际上，那些使 Whitehead 定理失效的“性质不良”空间在类型论中是不可见的。
粗略来说，性质良好的拓扑空间已经足以呈现 $\infty$-群胚，%
\index{.infinity-groupoid@$\infty$-groupoid}
而同伦类型论直接处理 $\infty$-群胚，而非具体的拓扑空间。
因此，人们或许会预期 Whitehead 定理在泛等基础中成立。

然而，情况并非如此：Whitehead 定理是不可证明的。
事实上，在已知的某些类型论模型中，该定理不成立，尽管其原因与在性质不良的拓扑空间上失效的原因完全不同。
这些模型是“非超完备 $\infty$-意象”
\index{.infinity1-topos@$(\infty,1)$-topos}%
\index{.infinity1-topos@$(\infty,1)$-topos!non-hypercomplete}%
（参见~\cite{lurie:higher-topoi}）；粗略地说，它们由 $\infty$ 维底空间上的 $\infty$-群胚层构成。

\index{axiom!Whitehead's principle|(}%
\index{Whitehead's!principle|(}%

因此，从基础的观点来看，我们可以将\emph{Whitehead 原理}视为一条“经典性公理”，类似于 \LEM{} 和 \choice{}。
它可以被一致地假设，但并不属于以计算为动机的类型论，也并非在所有自然模型中成立。
但在基于集合论的基础下，这条原理是隐而不显的：在一个 $\infty$-群胚由集合（使用拓扑空间、单纯集合等模型）构建起来的世界里，它不可能不成立。

这听起来或许有些奇怪，但实际上并不意外。
同伦类型论是同伦类型的\emph{抽象}理论，而在集合论框架下，拓扑空间或单纯集合的同伦论只是该理论的一个\emph{具体}模型；
这就像整数是抽象的环论的一个具体模型一样。
可以想见，任何具体模型都会具有一些并非相应抽象理论所固有的特殊性质，但有时我们希望将它们作为额外的公理来假定（例如，整数环是主理想整环，但并非所有环都是如此）。

描述类型论的模型已超出本书范围，因此我们不会解释 Whitehead 原理在某些模型中为何可能不成立。
然而，我们可以证明，只要所涉及的类型对某个有限 $n$ 是 $n$-截断的，该原理便成立，证明方法是对 $n$ 作“向下”归纳。
这一结果本身有其独立的意义（例如，它蕴含了艾伦伯格–麦克兰恩空间的本质唯一性）；
同时，其证明过程也有望提供一些直观解释，帮助我们理解为什么在缺少截断条件的情况下，我们无法指望证明类似的定理。

我们首先给出 \cref{thm:mono-surj-equiv} 的如下变形，它最终将为截断版 Whitehead 原理的归纳步骤提供支持。
它可以看作范畴论中经典结论“全忠实且本质满的函子是范畴等价”在类型论与 $\infty$-群胚语境下的版本。

\begin{thm}\label{thm:whitehead0}
  设 $f:A\to B$ 是一个函数，满足：
  \begin{enumerate}
  \item $\trunc0 f : \trunc0 A \to \trunc0 B$ 是满的，并且\label{item:whitehead01}
  \item 对任意 $x,y:A$，函数 $\apfunc f : (\id[A]xy) \to (\id[B]{f(x)}{f(y)})$ 是一个等价。\label{item:whitehead02}
  \end{enumerate}
  那么 $f$ 是一个等价。
\end{thm}

\begin{proof}
  注意条件~\ref{item:whitehead02} 恰好表达了 $f$ 是嵌入，参见~\cref{sec:mono-surj}。
  因此，由 \cref{thm:mono-surj-equiv}，只需证明 $f$ 是满的，即对任意 $b:B$ 证明 $\trunc{-1}{\hfib f b}$。
  给定 $b$。由于 $\trunc0 f$ 是满的，仅知存在 $a:A$ 使得 $\trunc 0 f(\tproj0a) = \tproj0b$。
  而我们的目标是命题，故可以假设给定了这样一个 $a$。
  于是有 $\tproj0{f(a)} = \trunc 0 f(\tproj0a) =\tproj0b$，从而 $\trunc{-1}{f(a)=b}$。
  同样，因为目标仍是命题，我们可以进一步假设 $f(a)=b$。
  由此 $\hfib f b$ 有实例，因此仅知有实例。
\end{proof}

由于同伦群是环路空间\index{loop space}的截断，而非道路空间的截断，我们需要将上述定理修改为关于环路空间的表述。
回忆从 \cref{def:loopfunctor} 引入的映射 $\Omega f$。

\begin{cor}\label{thm:whitehead1}
  设 $f:A\to B$ 是一个函数，满足：
  \begin{enumerate}
  \item $\trunc0 f : \trunc0 A \to \trunc0 B$ 是一个双射，并且
  \item 对任意 $x:A$，函数 $\Omega f : \Omega(A,x) \to \Omega(B,f(x))$ 是一个等价。
  \end{enumerate}
  那么 $f$ 是一个等价。
\end{cor}

\begin{proof}
  由 \cref{thm:whitehead0}，只需证明对任意 $x,y:A$，$\apfunc f : (\id[A]xy) \to (\id[B]{f(x)}{f(y)})$ 是一个等价。
  而根据 \cref{thm:equiv-inhabcod}，我们可以假设 $\id[B]{f(x)}{f(y)}$。
  特别地，有 $\tproj0{f(x)} = \tproj0{f(y)}$，因此由于 $\trunc0 f$ 是等价，可知 $\tproj0 x = \tproj0y$，故 $\tproj{-1}{x=y}$。
  由于我们所要证明的是命题（即 $\apfunc f$ 为等价），我们可以假设给定 $p:x=y$。
  但此时以下方块在同伦意义下交换：
  \begin{equation*}
  \vcenter{\xymatrix@C=3pc{
      \Omega(A,x)\ar[r]^-{\blank\ct p}\ar[d]_{\Omega f} &
      (\id[A]xy) \ar[d]^{\apfunc f}\\
      \Omega(B,f(x))\ar[r]_-{\blank\ct f(p)} &
      (\id[B]{f(x)}{f(y)}).
      }}
  \end{equation*}
  顶边和底边的映射都是等价，而左边的映射由假设也是等价。
  因此，根据“三中取二”性质，右边的映射也是等价。
\end{proof}

现在我们可以证明截断版 Whitehead 原理。

\begin{thm}\label{thm:whiteheadn}
  设 $A$ 与 $B$ 是 $n$-类型，且 $f:A\to B$ 满足：
  \begin{enumerate}
  \item $\trunc0f:\trunc0A \to \trunc0B$ 是一个双射，并且\label{item:wh0}
  \item $\pi_k(f):\pi_k(A,x) \to \pi_k(B,f(x))$ 对所有 $k\ge 1$ 和所有 $x:A$ 都是双射。\label{item:whk}
  \end{enumerate}
  那么 $f$ 是一个等价。
\end{thm}

\noindent
条件~\ref{item:wh0} 几乎就是~\ref{item:whk} 中 $k=0$ 的情形，只是它没有提及任何基点 $x:A$。

\begin{proof}
  我们通过对 $n$ 的归纳来证明。
  当 $n=-2$ 时，结论是平凡的。
  因此，假设结论对所有 $n$-类型之间的函数都成立，并令 $A$ 和 $B$ 为 $(n+1)$-类型，且 $f:A\to B$ 如上所述。
  \cref{thm:whitehead1} 中的第一个条件由假设成立，因此只需证明对于任意 $x:A$，函数 $\Omega f: \Omega(A,x) \to \Omega(B,f(x))$ 是一个等价。

  由于 $\Omega(A,x)$ 和 $\Omega(B,f(x))$ 是 $n$-类型，我们可以应用归纳假设。
  我们需要验证 $\trunc 0{\Omega f}$ 是双射，并且对于所有 $k\geq1$ 和 $p : x = x$，映射 $\pi_k(\Omega f):\pi_k(x = x,p) \to \pi_k(f(x) = f(x),\Omega f(p))$ 是双射。
  第一个论断由假设成立，因为 $\trunc 0{\Omega f} \jdeq \pi_1(f)$。
  为证明第二个论断，我们首先将其推广：我们证明对于所有 $y : A$ 和 $q : x = y$，有 $\pi_k(\apfunc f):\pi_k(x = y,q) \to \pi_k(f(x) = f(y),\apfunc f(q))$。
  这蕴含了待证结论，因为当 $y\defeq x$ 时，在等同其基点 $\Omega f(p)=\apfunc f(p)$ 的意义下，有 $\pi_k(\Omega f)=\pi_k(\apfunc f)$。
  为证明这一推广，由道路归纳，只需证 $q$ 为 $\refl{a}$ 的情形。
  此时有 $\pi_k(\apfunc f) = \pi_k(\Omega f) = \pi_{k+1}(f)$，而 $\pi_{k+1}(f)$ 由原始假设为双射。
\end{proof}

注意，如果 $A$ 和 $B$ 对任何有限的 $n$ 都不是 $n$-类型，那么归纳法就无从开始。

\begin{cor}\label{thm:whitehead-contr}
  如果 $A$ 是一个 $0$-连通的 $n$-类型，并且对于所有 $k$ 和 $a:A$ 有 $\pi_k(A,a)=0$，那么 $A$ 是可缩的。
\end{cor}

\begin{proof}
  将 \cref{thm:whiteheadn} 应用到映射 $A\to\unit$ 上。
\end{proof}

作为一个应用，我们可以推断出 \cref{thm:conn-pik} 的逆命题。

\begin{cor}\label{thm:pik-conn}
  对于 $n\ge 0$，一个映射 $f:A\to B$ 是 $n$-连通的，当且仅当下述所有条件成立：
  \begin{enumerate}
  \item $\trunc0f:\trunc0A \to \trunc0B$ 是一个同构。
  \item 对于任意 $a:A$ 和 $k\le n$，映射 $\pi_k(f):\pi_k(A,a) \to \pi_k(B,f(a))$ 是一个同构。
  \item 对于任意 $a:A$，映射 $\pi_{n+1}(f):\pi_{n+1}(A,a) \to \pi_{n+1}(B,f(a))$ 是满的。
  \end{enumerate}
\end{cor}

\begin{proof}
  “仅当”方向就是 \cref{thm:conn-pik}。
  反过来，由纤维化的长正合列（\cref{thm:les}），
  \index{sequence!exact}%
  这些前提蕴含了对于所有 $k\le n$ 和 $a:A$ 有 $\pi_k(\hfib f {f(a)})=0$，并且 $\trunc 0{\hfib f{f(a)}}$ 是可缩的。
  由于对于 $k\le n$ 有 $\pi_k(\hfib f {f(a)}) = \pi_k(\trunc n{\hfib f{f(a)}})$，并且 $\trunc n{\hfib f{f(a)}}$ 是 $n$-连通的，根据 \cref{thm:whitehead-contr}，对任意 $a$，它是可缩的。

  余下只需证明对于 $b:B$（不一定具有 $f(a)$ 的形式），$\trunc n{\hfib f{b}}$ 是可缩的。
  然而，由假设，存在 $x:\trunc0A$ 使得 $\tproj 0b = \trunc0f(x)$。
  因为可缩性是命题，我们可以假定 $x$ 具有 $\tproj0a$ 的形式（其中 $a:A$），在这种情况下 $\tproj 0 b = \trunc0f(\tproj0a) = \tproj0{f(a)}$，因此 $\trunc{-1}{b=f(a)}$。
  同样由于可缩性是命题，我们可以假定 $b=f(a)$，结论随之得出。
\end{proof}

一个使得 $\trunc0f$ 是双射且对所有 $k$ 有 $\pi_k(f)$ 是双射的映射 $f$，被称为 \define{$\infty$-连通的（$\infty$-connected）}%
\indexdef{function!.infinity-connected@$\infty$-connected}%
\indexdef{.infinity-connected function@$\infty$-connected function}
或一个\define{弱等价}。%
\indexdef{equivalence!weak}%
\indexdef{weak equivalence!of types}
这等价于要求 $f$ 对所有 $n$ 都是 $n$-连通的。
一个类型 $Z$ 被称为 \define{$\infty$-截断的（$\infty$-truncated）}%
\indexdef{type!.infinity-truncated@$\infty$-truncated}%
\indexdef{.infinity-truncated type@$\infty$-truncated type}
或\define{超完备的}%
\indexdef{type!hypercomplete}%
\indexdef{hypercomplete type}
如果只要 $f$ 是 $\infty$-连通的，诱导映射
\[(\blank\circ f):(B\to Z) \to (A\to Z)\]
就是一个等价——也就是说，如果 $Z$ 认为每个 $\infty$-连通映射都是等价。
\indexdef{axiom!Whitehead's principle}%
那么，如果我们想把 Whitehead 原理当作一条公理来假设，可以使用以下任意一种等价形式。

\begin{itemize}
\item 每个 $\infty$-连通函数都是一个等价。
\item 每个类型都是 $\infty$-截断的。
\end{itemize}

在高阶意象模型中，
\index{.infinity1-topos@$(\infty,1)$-topos}%
$\infty$-截断的类型按照 \cref{sec:modalities} 的意义构成一个反射子宇宙（即一个 $(\infty,1)$-意象的“超完备化”），但我们不知道这一性质在一般情况下是否成立。

\index{axiom!Whitehead's principle|)}%
\index{Whitehead's!principle|)}%

对任何 $n$ 都不是 $n$-类型的类型是否\emph{存在}，这或许并不显然，但事实上确实存在。
在经典同伦论中，对于任何 $n\ge 2$，$\Sn^n$ 都具有这一性质。
我们在同伦类型论中尚未证明这一事实，但存在其他类型，我们可以证明它们具有“无限截断层次”。

\begin{eg}
  假设我们有 $B:\nat\to\type$，使得对每个 $n$，类型 $B(n)$ 都包含一个与 $n$ 重自反性不相等的 $n$-环路\index{loop!n-@$n$-}，例如 $p_n:\Omega^n(B(n),b_n)$ 满足 $p_n \neq \refl{b_n}^n$。
  （举例来说，我们可以定义 $B(n)\defeq \Sn^n$，并令 $p_n$ 为 $1:\Z$ 在同构 $\pi_n(\Sn^n)\cong \Z$ 下的像。）
  考虑 $C\defeq \prd{n:\nat} B(n)$，其中的点 $c:C$ 由 $c(n)\defeq b_n$ 定义。
  因为环路空间与积交换，所以对任意 $m$，我们有
  \[\eqvspaced{\Omega^m (C,c)}{\prd{n:\nat}\Omega^m(B(n),b_n)}.\]
  在此等价下，$\refl{c}^m$ 对应于函数 $(n\mapsto \refl{b_n}^m)$。
  现在在右边的类型中定义 $q_m$ 为
  \[ q_m(n) \defeq
  \begin{cases}
    p_n &\quad m=n\\
    \refl{b_n}^m &\quad m\neq n.
  \end{cases}
  \]
  若 $q_m = (n\mapsto \refl{b_n}^m)$ 成立，则 $p_n = \refl{b_n}^n$ 也成立，但事实并非如此。
  因此，$q_m \neq (n\mapsto \refl{b_n}^m)$，于是在 $\Omega^m(C,c)$ 中存在一个与 $\refl{c}^m$ 不相等的点。
  故 $C$ 对任何 $m:\nat$ 都不是 $m$-类型。
\end{eg}

我们预期，利用宇宙包含所有高阶归纳类型（如对一切 $n$ 包含 $\Sn^n$）这一事实，也应该能够证明一个宇宙 $\UU$ 本身对任何 $n$ 都不是 $n$-类型。
不过，这一点尚未得到证明。

\section{编码-解码方法的一般陈述}
\label{sec:general-encode-decode}

\indexdef{encode-decode method}

我们已经使用编码-解码方法来刻画各种类型的道路空间，包括余积（\cref{thm:path-coprod}）、自然数（\cref{thm:path-nat}）、截断（\cref{thm:path-truncation}）、圆（\cref{{cor:omega-s1}}）、纬悬（\cref{thm:freudenthal}）和推出（\cref{thm:van-Kampen}）。
本章引言中提及的许多其他定理的证明也用到了这一方法的变体，例如直接证明 $\pi_n(\Sn^n)$、Blakers--Massey 定理以及构造艾伦伯格–麦克兰恩空间。
虽然很想将这一方法抽象为一个引理，但这很困难，因为不同的问题需要稍有不同的变体。
例如，同一方法的不同变体可以用来刻画一个环路空间（如 \cref{thm:path-coprod,cor:omega-s1}）或整个道路空间（如 \cref{thm:path-nat}）；可以给出环路空间的完整刻画（例如 $\Omega^1(\Sn ^1)$），也可以仅刻画其某个截断（例如 van Kampen 定理）；既可以用于计算同伦群，也可以用于证明一个映射是 $n$-连通的（例如 Freudenthal 定理和 Blakers--Massey 定理）。

然而，我们可以针对该方法的特定变体来陈述引理。
这些引理的证明几乎是平凡的；主要目的在于通过一般性地陈述来阐明这一方法。
最简单的情况是用编码-解码方法来刻画一个类型的环路空间，如 \cref{thm:path-coprod} 和 \cref{{cor:omega-s1}}。

\begin{lem}[环路空间的编码-解码]\label{lem:encode-decode-loop}
  \index{loop space}%
给定一个带点类型 $(A,a_0)$ 和一个纤维化 $\code : A \to \type$，如果
\begin{enumerate}
\item $c_0 : \code(a_0)$，\label{item:ed1}
\item $\decode : \prd{x:A} \code(x) \to (\id{a_0}{x})$，\label{item:ed2}
\item 对所有 $c : \code(a_0)$ 有 \id{\transfib{\code}{\decode(c)}{c_0}}{c}，并且\label{item:ed3}
\item $\id{\decode(c_0)}{\refl{}}$，\label{item:ed4}
\end{enumerate}
那么 $(\id{a_0}{a_0})$ 等价于 $\code(a_0)$。
\end{lem}

\begin{proof}
定义
$\encode : \prd{x:A} (\id{a_0}{x}) \to \code(x)$ 为
\[
\encode_x(\alpha) = \transfib{\code}{\alpha}{c_0}.
\]
我们证明 $\encode_{a_0}$ 和 $\decode_{a_0}$ 互为拟逆。
复合 $\encode_{a_0} \circ \decode_{a_0}$ 由假设~\ref{item:ed3} 立即可得。
对于另一个复合，我们证明
\[
\prd{x:A}{p : \id{a_0}{x}} \id{\decode_{x} (\encode_{x} p)}{p}.
\]
由道路归纳，只需证明
$\id{{\decode_{{a_0}} (\encode_{{a_o}} \refl{})}}{\refl{}}$。
在归约 $\mathsf{transport}$ 之后，只需证明
$\id{{\decode_{{a_0}} (c_0)}}{\refl{}}$，而这正是假设~\ref{item:ed4}。
\end{proof}

如果想要得到 $(\id{a_0}{\blank})$ 与 $\code$ 之间的逐纤维等价，只需将条件 (iii) 加强为
\[
\prd{x:A}{c : \code(x)} \id{\encode_{x}(\decode_{x}(c))}{c}.
\]
然而，为了计算一个环路空间（例如 $\Omega(\Sn ^1)$），这一更强的假设并非必要。

另一个常见的变体出现在计算同伦群时，它刻画的是环路空间的截断：

\begin{lem}[环路空间截断的编码-解码]
假设有一个带点类型 $(A,a_0)$ 和一个纤维化 $\code : A \to \type$，其中对每个 $x$，$\code(x)$ 都是 $k$-类型。
定义
\[
\encode : \prd{x:A} \trunc{k}{\id{a_0}{x}} \to \code(x)
\]
（利用 $\code(x)$ 是 $k$-类型这一事实）通过截断递归，将 $\alpha : \id{a_0}{x}$ 映到 $\transfib{\code}{\alpha}{c_0}$。假设：
\begin{enumerate}
\item $c_0 : \code(a_0)$，
\item $\decode : \prd{x:A} \code(x) \to \trunc{k}{\id{a_0}{x}}$，
\item \label{item:decode-encode-loop-iii}
  $\id{\encode_{a_0}(\decode_{a_0}(c))}{c}$ 对所有 $c : \code(a_0)$ 成立，并且
\item \label{item:decode-encode-loop-iv}
  $\id{\decode(c_0)}{\tproj{}{\refl{}}}$。
\end{enumerate}
那么 $\trunc{k}{\id{a_0}{a_0}}$ 等价于 $\code(a_0)$。
\end{lem}

\begin{proof}
$\decode \circ \encode$ 为恒等映射，这由 \ref{item:decode-encode-loop-iii} 立即可得。
%
为证明 $\encode \circ \decode$，我们首先进行截断归纳，由此只需证明
\[
\prd{x:A}{p : \id{a_0}{x}} \id{\decode_{x}(\encode_{x}(\tproj{k}{p}))}{\tproj{k}{p}}.
\]
之所以可以进行截断归纳，是因为 $k$-类型中的道路也是 $k$-类型。
为证明此类型，我们进行道路归纳，并在归约 \encode 之后，使用假设~\ref{item:decode-encode-loop-iv}。
\end{proof}

\section{更多结果}
\label{sec:moreresults}

尽管本章没有给出证明，但下列结果也已在同伦类型论中建立。

\begin{thm}
\index{homotopy!group!of sphere}
存在一个 $k$，使得对所有 $n \ge 3$，$\pi_{n+1}(\Sn ^n) =
\Z_k$。
\end{thm}

\begin{proof}[关于证明的注记.]
证明包括对 $\pi_4(\Sn ^3)$ 的计算，以及诉诸稳定性结论（\cref{cor:stability-spheres}）。
在这一结果的经典陈述中，$k$ 等于 $2$。
虽然我们尚未验证 $k$ 确实为 $2$，但我们对 $\pi_4(\Sn ^3)$ 的计算是构造性的\index{mathematics!constructive}，与本章中的所有其他证明一样。
（更准确地说，它没有使用任何额外公理，例如 \LEM{} 或 \choice{}，从而使其构造性与泛等性和高阶归纳类型相当。）
因此，给定同伦类型论的一种计算性解释，我们就能够在计算机上运行该证明，从而验证 $k$ 确实为 $2$。
这个例子颇耐人寻味，因为这是首个无需事先知道答案即可完成的同伦群计算。
\end{proof}

% Recall from \cref{sec:colimits} that $X \sqcup^C Y$ denotes the
% (homotopy) pushout of $X$ and $Y$ along $C$.
\index{pushout}%

\begin{thm}[Blakers--Massey 定理]\label{Blakers-Massey}
  \indexdef{theorem!Blakers--Massey}%
  \indexsee{Blakers--Massey theorem}{theorem, Blakers--Massey}%
  假设给定映射 $f : C  \rightarrow X$ 和 $g : C \rightarrow Y$。
  先取 $f$ 与 $g$ 的推出 $X \sqcup^C Y $，再取其内含映射的拉回 $\inl : X \rightarrow X \sqcup^C Y \leftarrow Y : \inr$，便得到一个诱导映射 $C \to X \times_{(X \sqcup^C Y)} Y$。

  如果 $f$ 是 $i$-连通的而 $g$ 是 $j$-连通的，那么这个诱导映射就是 $(i+j)$-连通的。
  换句话说，对于任意点 $x:X$ 和 $y:Y$，对应的纤维 $C_{x,y}$（即 $(f,g) : C \to X \times Y $ 的纤维）给出了推出中道路空间 $\id[X \sqcup^C Y]{\inl(x)}{\inr(y)}$ 的一个近似。
\end{thm}

需要注意的是，在经典代数拓扑中，Blakers--Massey 定理常常以稍有不同的形式陈述：
映射 $f$ 和 $g$ 被替换为 CW 复形子复形的包含映射，而同伦推出与同伦拉回则分别被替换为并集与交集。
为了在同伦类型论中表达这一定理，我们必须用同伦不变的概念来替换此类概念。
在 van Kampen 定理（\cref{sec:van-kampen}）中我们已经看到了另一个类似的例子：在那里，我们必须用同伦推出来替换开子集的并集。

\begin{thm}[Eilenberg--Mac Lane 空间]\label{Eilenberg-Mac-Lane-Spaces}
\index{Eilenberg--Mac Lane space}
对于任意交换\index{group!abelian}群 $G$ 和正整数 $n$，都存在一个 $n$-类型
$K(G,n)$，使得 $\pi_n(K(G,n)) = G$，并且当 $k\neq n$ 时，$\pi_k(K(G,n)) = 0$。
\end{thm}

\begin{thm}[覆叠空间]\label{thm:covering-spaces}
  \index{covering space}%
  对于连通空间 $A$，$A$ 上的覆叠空间与带有 $\pi_1(A)$ 作用的集合之间存在等价。
\end{thm}

\sectionNotes

%% For the spheres,  two
%% different definitions of the $n$-sphere $\Sn ^n$ have been used: the first as the
%% suspension of $\Sn ^ {n-1}$ (\cref{sec:suspension}), and the second
%% as a higher inductive type with one base point and one loop in
%% $\Omega^n$ (\cref{sec:circle}); we list the status for both
%% definitions.  The equivalence of the two can be deduced from the remarks
%% at the end of \cref{sec:suspension}.

{
\newcommand{\humancheck}{\ding{52}}
\newcommand{\computercheck}{\ding{52}\kern-0.5em\ding{52}}


\begin{table}[htb]
  \centering
\begin{tabular}{lcc}
\toprule
定理         & 状态 \\
\midrule
$\pi_1(\Sn ^1)$                     & \computercheck \\
$\pi_{k<n}(\Sn ^n)$                  & \computercheck \\
%% $\pi_{k<n}(\Sn ^n)$ --- suspension  & \computercheck \\
%% $\pi_2(\Sn ^2)$ --- suspension      & \computercheck \\
同伦群的长正合列 & \computercheck    \\
Hopf 纤维化的全空间是 $\Sn ^3$ & \humancheck    \\
$\pi_2(\Sn ^2)$                     & \computercheck \\
$\pi_3(\Sn ^2)$                     & \humancheck    \\
$\pi_n(\Sn ^n)$                     & \computercheck \\
%% $\pi_n(\Sn ^n)$ --- suspension      & \computercheck \\
$\pi_4(\Sn ^3)$                     & \humancheck    \\
Freudenthal 纬悬定理      & \computercheck \\
Blakers--Massey 定理              & \computercheck \\
Eilenberg--Mac Lane 空间 $K(G,n)$ & \computercheck \\
van Kampen 定理                & \computercheck \\
覆叠空间                     & \computercheck \\
关于 $n$-类型的 Whitehead 原理 & \computercheck \\
\bottomrule
\end{tabular}
\caption{同伦论中定理的手工证明（\humancheck）与计算机证明（\computercheck）情况。}
  \label{tab:theorems}
\end{table}


}

本章所描述的定理都是经典同伦论中的标准结果；其中许多在 \cite{hatcher02topology} 中有所记述。
在这部分附注里，我们回顾这些定理在同伦类型论中新的综合证明是如何发展起来的。
\cref{tab:theorems} 列出了已在同伦类型论中得到证明的同伦论定理，以及它们是否经过了计算机验证。
这些成果几乎全部诞生于 2013 年高等研究院（IAS）的春季学期，是一次重要协作努力的结晶。
许多人为此作出了贡献，例如担任某个证明的主要作者、提出值得研究的问题、参与大量相关问题的讨论与研讨班，或对已有成果提供反馈。
以下人员是上述定理的首批同伦类型论证明的主要作者。
除非另有说明，对于已经过计算机验证的定理，其主要作者也是率先使用计算机证明助理将证明\index{mathematics!formalized}形式化的人。

\begin{itemize}
\item
Shulman 给出了 $\pi_1(\Sn^1)$ 的同伦论计算。Licata 后来发现了编码-解码证明以及编码-解码方法。

\item
Brunerie 计算了 $\pi_{k<n}(\Sn ^ n)$。
Licata 后来给出了一个编码-解码的版本。

\item Voevodsky（沃沃茨基）构造了同伦群的长正合列。

\item Lumsdaine 构造了 Hopf 纤维化。
  Brunerie 证明了其全空间是 $\Sn ^3$，从而算出了 $\pi_2(\Sn ^2)$ 和 $\pi_3(\Sn ^3)$。

\item Licata 和 Brunerie 直接计算了 $\pi_n(\Sn ^n)$。

\item
  Lumsdaine 证明了 Freudenthal 纬悬定理；Licata 和 Lumsdaine 将这一证明形式化。
\item Lumsdaine, Finster 和 Licata 证明了 Blakers--Massey 定理；
  Lumsdaine, Brunerie, Licata 和 Hou 将其形式化。

\item
Licata 使用编码-解码方法计算了 $\pi_2(\Sn ^2)$，并利用 Freudenthal 纬悬定理计算了 $\pi_n(\Sn ^n)$；运用类似技术，他构造了 $K(G,n)$。

\item
Shulman 证明了 van Kampen 定理；Hou 将此证明形式化。

\item
Licata 证明了关于 $n$-类型的 Whitehead 定理。

\item Brunerie 计算了 $\pi_4(\Sn ^3)$。

\item
Hou 建立并形式化了覆叠空间的理论。
\end{itemize}

同伦论与类型论之间的相互交融对于这些成果的发展至关重要。
例如，$\pi_1(\Sn ^1)=\mathbb{Z}$ 的第一个证明是 \cref{subsec:pi1s1-homotopy-theory} 中给出的那个，它沿用了经典同伦论的思路。
对这一证明的类型论分析催生了编码-解码方法。
对 $\pi_2(\Sn ^2)$ 的首次计算也沿用了经典方法，但这很快便引出了该结果的一个编码-解码证明。
编码-解码计算被推广到 $\pi_n(\Sn ^n)$，这进而引出了 Freudenthal 纬悬定理的证明——其方法是将编码-解码论证与关于连通性的经典同伦论推理相结合。
这一证明反过来又催生了 Blakers--Massey 定理与 Eilenberg--Mac Lane 空间。
这一系列成果的快速进展，彰显了我们对这两门学科之间联系的新理解所蕴含的潜力。

\sectionExercises

\begin{ex}\label{ex:homotopy-groups-resp-prod}
  证明同伦群保持积：$\eqv{\pi_n(A\times B)}{\pi_n(A)\times \pi_n(B)}$。
\end{ex}

\begin{ex}\label{ex:decidable-equality-susp}
  \index{decidable!equality}%
  证明：若 $A$ 是一个具备可判定相等性（参见 \cref{defn:decidable-equality}）的集合，则其纬悬 $\susp A$ 是一个 1-类型。
  （这是否可以不借助可判定相等性的假设来证明，仍是一个\index{open!problem}开放问题。）
\end{ex}

\begin{ex}\label{ex:contr-infinity-sphere-colim}
  定义 $\Sn^\infty$ 为序列 $\Sn^0 \to \Sn^1 \to \Sn^2 \to\cdots$ 的余极限。
  证明 $\Sn^\infty$ 是可缩的。
\end{ex}

\begin{ex}\label{ex:contr-infinity-sphere-susp}
  将 $\Sn^\infty$ 定义为由以下生成元生成的高阶归纳类型：
  \begin{itemize}
  \item 两个点 $\north:\Sn^\infty$ 和 $\south:\Sn^\infty$，及
  \item 对每个 $x:\Sn^\infty$，一条道路 $\merid(x):\north=\south$。
  \end{itemize}
  换言之，$\Sn^\infty$ 就是它自身的纬悬。
  证明 $\Sn^\infty$ 是可缩的。
\end{ex}

\begin{ex}\label{ex:unique-fiber}
  设 $f:X\to Y$ 是一个映射且 $Y$ 是连通的。
  证明对任意 $y_1,y_2:Y$，有 $\brck{\eqv{\hfib f {y_1}}{\hfib f {y_2}}}$。
\end{ex}

\begin{ex}\label{ex:ap-path-inversion}
  对任意带点类型 $A$，令 $i_A : \Omega A \to \Omega A$ 表示环路取逆，即 $i_A \defeq \lam{p} \rev{p}$。
  证明 $i_{\Omega A} : \Omega^2 A \to \Omega^2 A$ 等于 $\Omega(i_A)$。
\end{ex}

\begin{ex}\label{ex:pointed-equivalences}
  将\textbf{带点等价}定义为其底层函数为等价的带点映射。
  \begin{enumerate}
  \item 证明带点类型 $(X,x_0)$ 与 $(Y,y_0)$ 之间的带点等价类型等价于 $\id[\pointed\type]{(X,x_0)}{(Y,y_0)}$。
  \item 用带点拟逆和带点同伦，按照 \cref{cha:equivalences} 中的某种相干风格，重新表述带点等价的概念。
  \end{enumerate}
\end{ex}

\begin{ex}\label{ex:HopfJr}
  仿照 \cref{sec:hopf} 中 Hopf 纤维化的例子，将 \define{低阶 Hopf 纤维化}
  \indexdef{Hopf!fibration!junior}%
定义为 $\Sn ^1$ 上的纤维化（即类型族），其在基点处的纤维为 $\Sn ^0$，全空间为 $\Sn ^1$。这也称为圆 $\Sn ^1$ 的``扭曲二重覆叠''。
\end{ex}

\begin{ex}\label{ex:SuperHopf}
再次仿照 \cref{sec:hopf} 中 Hopf 纤维化的例子，定义 $\Sn ^4$ 上的一个类似纤维化，其在基点处的纤维为 $\Sn ^3$，全空间为 $\Sn ^7$。这是同伦类型论中的一个\index{open!problem}未解决问题（在经典同伦论中已知存在这样的纤维化）。
\end{ex}

\begin{ex}\label{ex:vksusppt}
  继续 \cref{eg:suspension} 的讨论，证明若 $A$ 有一个点 $a:A$，则可将 $\pi_1(\susp A)$ 等同于以 $\trunc0A$ 为生成元、以关系 $\tproj0a = e$ 给出的群。
  进而证明，若假设排中律，该群也是 $\trunc0 A \setminus \{\tproj0 a \}$ 上的自由群。
\end{ex}

\begin{ex}\label{ex:vksuspnopt}
  再次继续 \cref{eg:suspension} 的讨论，但这次不假定 $A$ 带点，证明可将 $\pi_1(\susp A)$ 等同于以生成元 $\trunc0A \times \trunc0A$ 和关系
  \begin{equation*}
    (a,b) = \opp{(b,a)},
    \qquad
    (a,c) = (a,b)\cdot (b,c),
    \qquad\text{and}\qquad
    (a,a) = e.
  \end{equation*}
  给出的群。
\end{ex}

\index{acceptance|)}

%%% Local Variables:
%%% mode: latex
%%% TeX-master: "hott-online"
%%% End: